\documentclass[fleqn,10pt]{wlscirep}
%DIF LATEXDIFF DIFFERENCE FILE
%DIF DEL paperV2/scientific_reports/main2.tex   Thu Aug 20 10:48:49 2026
%DIF ADD paperV3/scientific_reports/main3.tex   Fri Sep  4 21:43:21 2026
\usepackage[utf8]{inputenc}
\usepackage[T1]{fontenc}
\usepackage{amsmath,amssymb}
\usepackage{array,booktabs,tabularx}
\usepackage{graphicx}
\usepackage{float}
\usepackage{algorithm}
\usepackage{algpseudocode}
\usepackage{lineno}
\usepackage{microtype}
\usepackage{url}
\usepackage{seqsplit}

\newcolumntype{Y}{>{\raggedright\arraybackslash}X}
\newcolumntype{L}[1]{>{\raggedright\arraybackslash}p{#1}}
\DeclareRobustCommand{\code}[1]{\path{#1}}
\newcommand{\Rank}{\operatorname{rank}}
\newcommand{\tok}{\operatorname{token}}
\newcommand{\Enc}{\operatorname{Enc}}
\newcommand{\Dec}{\operatorname{Dec}}
\newcommand{\RankOf}{\operatorname{rank\_of}}
\newcommand{\Encode}{\operatorname{encode}}
\newcommand{\Decode}{\operatorname{decode}}
\newcommand{\Tokenize}{\operatorname{tokenize}}
\newcommand{\Detokenize}{\operatorname{detokenize}}
\newcommand{\NextLogits}{\operatorname{next\_logits}}
\modulolinenumbers[5]

\title{\DIFdel{RankCloak: Hiding Synthetic Cryptographic Artifacts in Plain English with Language Models}\\[0.6em]\DIFadd{RankCloak Conceals the Surface Form of Synthetic Cryptographic Artifacts in Language Model Generated Text}}
%DIF 31a31
\hypersetup{pdftitle={RankCloak Conceals the Surface Form of Synthetic Cryptographic Artifacts in Language Model Generated Text},pdfauthor={Alexander V. Mantzaris}} %DIF >
%DIF -------

\author[1,*]{Alexander V. Mantzaris}
\affil[1]{School of Data, Mathematical, and Statistical Sciences, University of Central Florida, Orlando, FL, United States of America}
\affil[*]{Correspondence: alexander.mantzaris@ucf.edu}

\keywords{linguistic steganography, large language models, rank transcoding, exact-copy replay, neural steganalysis}

\begin{abstract}
\textbf{Deleted V2 abstract}\par
\DIFdelbegin\DIFdel{RankCloak transcodes synthetic cryptographic artifacts into natural language using language-model next-token ranks (leveraging an LLM) and realizes those ranks as generated cover text. We compare direct subword representation with bounded ASCII-byte and hex-nibble codecs, including non-segmented generation and segmented forced spans followed by non-payload natural tails. The completed evaluation used 480 public deterministic artifacts from eight cryptographic classes, three open-source model families, and 18 English prompt templates. The secondary recovery tests used Spanish and Simplified Chinese prompts. All 6,480 primary payload trials recovered the serialized payload under saved-token exact-copy replay (Wilson 95\% confidence interval \([0.999408, 1.000000]\)), and all 576 multilingual trials recovered under the same condition. This mapping contract is restrictive. Detokenizing and retokenizing visible text recovered 61.1\%; normalization and formatting changes reduced recovery further. Using markdown transfer, paraphrase, and cross-model decoding recovered none. Automated readability diagnostics were protocol-dependent and are not participant ratings. Neural detectors strongly distinguished RankCloak covers from the controls (matched ROC-AUC 0.990190 for TextCNN and 0.999834 for DeBERTa-v3-base), providing adverse evidence for concealment. RankCloak therefore provides a reproducible framework for measuring rank pressure, capacity, replay requirements, cover cost, and detectability. It does not establish a secure or deployment-ready covert channel.}\DIFdelend
\par\medskip
\textbf{Added final V3 abstract}\par
\DIFaddbegin\DIFadd{Cryptographic artifacts such as hashes, nonces, identifiers, and ciphertexts are conspicuous when transmitted in their usual textual forms. RankCloak conceals this recognizable form by encoding each synthetic artifact through language-model token choices, producing cover text that does not directly display the artifact. We evaluated 480 public deterministic artifacts from eight classes with three open-source model families and 18 English prompt templates. The decoder recovered every artifact across 6,480 primary trials when the original token sequence and shared generation configuration were preserved. Recovery from rendered text was 61.1 percent; markdown transfer, paraphrasing, and decoding with another model failed in the tested robustness conditions. Stronger detection tests separated duplicate and near-duplicate material before training, withheld model families, included human-written comparisons, and used generating-model probabilities; encoded covers were often detectable. Restricting embedding to positions with several plausible next words reduced some detection signals but also reduced capacity and increased cover length. Detecting a possible encoding does not itself display or reconstruct the concealed artifact. RankCloak demonstrates artifact-specific concealment and recovery under a shared configuration while exposing detectability and transmission limitations; it does not establish hiding arbitrary natural-language messages.}\DIFaddend
\end{abstract}
%DIF PREAMBLE EXTENSION ADDED BY LATEXDIFF
%DIF UNDERLINE PREAMBLE %DIF PREAMBLE
\RequirePackage[normalem]{ulem} %DIF PREAMBLE
\RequirePackage{color}\definecolor{RED}{rgb}{1,0,0}\definecolor{BLUE}{rgb}{0,0,1} %DIF PREAMBLE
\providecommand{\DIFadd}[1]{{\protect\color{blue}\uwave{#1}}} %DIF PREAMBLE
\providecommand{\DIFdel}[1]{{\protect\color{red}\sout{#1}}}                      %DIF PREAMBLE
%DIF SAFE PREAMBLE %DIF PREAMBLE
\providecommand{\DIFaddbegin}{} %DIF PREAMBLE
\providecommand{\DIFaddend}{} %DIF PREAMBLE
\providecommand{\DIFdelbegin}{} %DIF PREAMBLE
\providecommand{\DIFdelend}{} %DIF PREAMBLE
\providecommand{\DIFmodbegin}{} %DIF PREAMBLE
\providecommand{\DIFmodend}{} %DIF PREAMBLE
%DIF FLOATSAFE PREAMBLE %DIF PREAMBLE
\providecommand{\DIFaddFL}[1]{\DIFadd{#1}} %DIF PREAMBLE
\providecommand{\DIFdelFL}[1]{\DIFdel{#1}} %DIF PREAMBLE
\providecommand{\DIFaddbeginFL}{} %DIF PREAMBLE
\providecommand{\DIFaddendFL}{} %DIF PREAMBLE
\providecommand{\DIFdelbeginFL}{} %DIF PREAMBLE
\providecommand{\DIFdelendFL}{} %DIF PREAMBLE
\newcommand{\DIFscaledelfig}{0.5}
%DIF HIGHLIGHTGRAPHICS PREAMBLE %DIF PREAMBLE
\RequirePackage{settobox} %DIF PREAMBLE
\RequirePackage{letltxmacro} %DIF PREAMBLE
\newsavebox{\DIFdelgraphicsbox} %DIF PREAMBLE
\newlength{\DIFdelgraphicswidth} %DIF PREAMBLE
\newlength{\DIFdelgraphicsheight} %DIF PREAMBLE
% store original definition of \includegraphics %DIF PREAMBLE
\LetLtxMacro{\DIFOincludegraphics}{\includegraphics} %DIF PREAMBLE
\newcommand{\DIFaddincludegraphics}[2][]{{\color{blue}\fbox{\DIFOincludegraphics[#1]{#2}}}} %DIF PREAMBLE
\newcommand{\DIFdelincludegraphics}[2][]{% %DIF PREAMBLE
\sbox{\DIFdelgraphicsbox}{\DIFOincludegraphics[#1]{#2}}% %DIF PREAMBLE
\settoboxwidth{\DIFdelgraphicswidth}{\DIFdelgraphicsbox} %DIF PREAMBLE
\settoboxtotalheight{\DIFdelgraphicsheight}{\DIFdelgraphicsbox} %DIF PREAMBLE
\scalebox{\DIFscaledelfig}{% %DIF PREAMBLE
\parbox[b]{\DIFdelgraphicswidth}{\usebox{\DIFdelgraphicsbox}\\[-\baselineskip] \rule{\DIFdelgraphicswidth}{0em}}\llap{\resizebox{\DIFdelgraphicswidth}{\DIFdelgraphicsheight}{% %DIF PREAMBLE
\setlength{\unitlength}{\DIFdelgraphicswidth}% %DIF PREAMBLE
\begin{picture}(1,1)% %DIF PREAMBLE
\thicklines\linethickness{2pt} %DIF PREAMBLE
{\color[rgb]{1,0,0}\put(0,0){\framebox(1,1){}}}% %DIF PREAMBLE
{\color[rgb]{1,0,0}\put(0,0){\line( 1,1){1}}}% %DIF PREAMBLE
{\color[rgb]{1,0,0}\put(0,1){\line(1,-1){1}}}% %DIF PREAMBLE
\end{picture}% %DIF PREAMBLE
}\hspace*{3pt}}} %DIF PREAMBLE
} %DIF PREAMBLE
\LetLtxMacro{\DIFOaddbegin}{\DIFaddbegin} %DIF PREAMBLE
\LetLtxMacro{\DIFOaddend}{\DIFaddend} %DIF PREAMBLE
\LetLtxMacro{\DIFOdelbegin}{\DIFdelbegin} %DIF PREAMBLE
\LetLtxMacro{\DIFOdelend}{\DIFdelend} %DIF PREAMBLE
\DeclareRobustCommand{\DIFaddbegin}{\DIFOaddbegin \let\includegraphics\DIFaddincludegraphics} %DIF PREAMBLE
\DeclareRobustCommand{\DIFaddend}{\DIFOaddend \let\includegraphics\DIFOincludegraphics} %DIF PREAMBLE
\DeclareRobustCommand{\DIFdelbegin}{\DIFOdelbegin \let\includegraphics\DIFdelincludegraphics} %DIF PREAMBLE
\DeclareRobustCommand{\DIFdelend}{\DIFOaddend \let\includegraphics\DIFOincludegraphics} %DIF PREAMBLE
\LetLtxMacro{\DIFOaddbeginFL}{\DIFaddbeginFL} %DIF PREAMBLE
\LetLtxMacro{\DIFOaddendFL}{\DIFaddendFL} %DIF PREAMBLE
\LetLtxMacro{\DIFOdelbeginFL}{\DIFdelbeginFL} %DIF PREAMBLE
\LetLtxMacro{\DIFOdelendFL}{\DIFdelendFL} %DIF PREAMBLE
\DeclareRobustCommand{\DIFaddbeginFL}{\DIFOaddbeginFL \let\includegraphics\DIFaddincludegraphics} %DIF PREAMBLE
\DeclareRobustCommand{\DIFaddendFL}{\DIFOaddendFL \let\includegraphics\DIFOincludegraphics} %DIF PREAMBLE
\DeclareRobustCommand{\DIFdelbeginFL}{\DIFOdelbeginFL \let\includegraphics\DIFdelincludegraphics} %DIF PREAMBLE
\DeclareRobustCommand{\DIFdelendFL}{\DIFOaddendFL \let\includegraphics\DIFOincludegraphics} %DIF PREAMBLE
%DIF COLORLISTINGS PREAMBLE %DIF PREAMBLE
\RequirePackage{listings} %DIF PREAMBLE
\RequirePackage{color} %DIF PREAMBLE
\lstdefinelanguage{DIFcode}{ %DIF PREAMBLE
%DIF DIFCODE_UNDERLINE %DIF PREAMBLE
  moredelim=[il][\color{red}\sout]{\%DIF\ <\ }, %DIF PREAMBLE
  moredelim=[il][\color{blue}\uwave]{\%DIF\ >\ } %DIF PREAMBLE
} %DIF PREAMBLE
\lstdefinestyle{DIFverbatimstyle}{ %DIF PREAMBLE
	language=DIFcode, %DIF PREAMBLE
	basicstyle=\ttfamily, %DIF PREAMBLE
	columns=fullflexible, %DIF PREAMBLE
	keepspaces=true %DIF PREAMBLE
} %DIF PREAMBLE
\lstnewenvironment{DIFverbatim}{\lstset{style=DIFverbatimstyle}}{} %DIF PREAMBLE
\lstnewenvironment{DIFverbatim*}{\lstset{style=DIFverbatimstyle,showspaces=true}}{} %DIF PREAMBLE
%DIF END PREAMBLE EXTENSION ADDED BY LATEXDIFF

\begin{document}
\setlength{\emergencystretch}{3em}
\sloppy
\maketitle
\linenumbers
\raggedbottom

\section*{Introduction}\label{sec:introduction}

Generative language models provide a new medium for text steganography because a prompt can define a next-token distribution and a topic direction in which generated text appears. In the Calgacus rank-transcoding construction\cite{NorelliBronstein2026Calgacus}, a sender tokenizes an arbitrary hidden text, records each token's context-dependent rank, and generates a cover text by selecting the correspondingly ranked tokens under a shared cover prompt; the receiver reverses these steps under the same contexts. Calgacus thereby transcodes one hidden-text token into one cover-text token. RankCloak retains that rank-transfer mechanism but also maps exact serialized cryptographic artifacts through direct-subword and bounded codecs, for which payload and cover token counts need not equal the Calgacus baseline. This mechanism connects classical work on covert channels, information-theoretic steganography, information hiding, and provably secure steganography\cite{Simmons1984Prisoners,Cachin1998InformationTheoretic,Hopper2002ProvablySecureSteganography,Petitcolas1999InformationHidingSurvey} with neural and LLM-based linguistic steganography, where hidden information is embedded through choices made during language generation\cite{Fang2017GeneratingSteganographicTextLSTMs,Yang2019RNNStega,Ziegler2019NeuralLinguisticSteganography,DaiCai2019NearImperceptible,Shen2020SelfAdjustingArithmetic,Zhang2021ProvablySecureGenerative,Wang2025DAIRStega,Wu2024GenerativeTextSteganographyLLM,Lin2024ZeroShotGLS}. Related work on LLM watermarking also treats token selection as a carrier of hidden signals, although watermarking \DIFdelbegin \DIFdel{is usually aimed at attribution or detection rather than arbitrary payload recovery\mbox{%DIFAUXCMD
\cite{Kirchenbauer2023WatermarkLLM}}\hskip0pt%DIFAUXCMD
. }\DIFdelend \DIFaddbegin \DIFadd{generally asks whether text carries a detectable mark rather than whether a chosen payload can be recovered\mbox{%DIFAUXCMD
\cite{Kirchenbauer2023WatermarkLLM}}\hskip0pt%DIFAUXCMD
. Entropy-guided watermarking selects suitable high-entropy positions for test-time marking\mbox{%DIFAUXCMD
\cite{Cai2025EntropyGuidedWatermarking}}\hskip0pt%DIFAUXCMD
; that idea motivates the selective-embedding experiment here, but RankCloak instead evaluates artifact-specific payload recovery, embedding rate, completion, cover length, and steganalysis.
}\DIFaddend

A long-running goal in this area is to make arbitrary, encrypted, or otherwise high-entropy data appear as ordinary communication. Early mimic-function work used statistical source models and inverse compression ideas to generate outputs with distributional properties similar to a target source\cite{Wayner1992MimicFunctions,Wayner1995StrongTheoreticalSteganography}. Hiding the Hidden and related linguistic-steganography systems transformed ciphertext into innocuous natural-language text that could later be transformed back into the original ciphertext\cite{ChapmanDavida1997HidingHidden,Bennett2004LinguisticSteganography}. Other text methods carried bits through synonym choice or contextual word substitution\cite{ChangClark2014PracticalLinguisticSteganography}, while network-camouflage systems generalized a similar objective to protocol formats by making encrypted traffic resemble allowed or innocuous traffic\cite{Weinberg2012StegoTorus,Dyer2013FormatTransformingEncryption}. RankCloak follows this same broad goal in an LLM setting. Instead of using a fixed grammar, synonym table, or protocol template, it uses the token choices made by a shared language model. The payload is first converted into ranks, and the cover prompt determines which English tokens occupy those ranks at each step. This also makes the method useful as a measurement tool since it exposes how a hidden artifact changes rank pressure, generated text quality, and the traces available for later detection.

This measurement focus also has defensive value for detecting the presence of such hidden encodings. If hashes, nonces, UUID-like values, or ciphertext strings are carried by generated tokens rather than displayed in their original syntax, surface scanners cannot rely only on visible hexadecimal, base64, UUID, or credential-shaped patterns. Detection must therefore look for indirect evidence, such as low-probability forced tokens, tokenizer artifacts, formatting fragments, unusually regular token choices, segmentation patterns, or differences between payload-bearing spans and non-payload continuations. Covert-channel and stegomalware surveys emphasize that hiding the existence of communication creates detection, capacity-limitation, and countermeasure problems across text and network media\cite{Zander2007CovertChannelsCountermeasures,Badar2025StegomalwareSurvey}. Text steganalysis and generated-text detection work similarly shows that statistical, neural, graph-based, and probability-based features can reveal some generated or steganographic text, while also showing that robust detection is setting-dependent and vulnerable to distribution shift or paraphrase\cite{Wen2019CNNTextSteganalysis,Wu2021GNNLinguisticSteganalysis,Wang2023FewShotLinguisticSteganalysis,Gehrmann2019GLTR,Mitchell2023DetectGPT,Sadasivan2025AIDetection}. The same issue is relevant to AI-safety monitoring, where hidden message passing and secret collusion among agents have been studied as ways that generated text might evade oversight\cite{Motwani2024SecretCollusion,Zolkowski2026EarlySigns}. RankCloak \DIFdelbegin \DIFdel{therefore treats hiding and detection as two sides of the same measurement problem. A useful methodology should recover payloads under declared exact-copy conditions while also producing transparent features, controls, and failure modes for steganalysis}\DIFdelend \DIFaddbegin \DIFadd{evaluates these as distinct endpoints. Artifact concealment asks whether the cryptographic artifact and its recognizable syntax are directly visible in the cover text. Payload recovery asks whether an authorized receiver using the shared recovery configuration can reconstruct that artifact, whereas steganalysis asks whether an observer can classify the cover as likely encoded. A positive steganalysis decision does not itself display or recover the concealed artifact, although high detectability weakens covert-channel stealth}\DIFaddend .

RankCloak proposes a methodology for a specific case, hiding synthetic cryptographic artifacts in generated English text. The payload classes resemble values commonly encountered in security, including hashes, authentication tags, nonces, identifiers, ciphertexts, and signatures\cite{NIST2015SecureHashStandard,Josefsson2006RFC4648,DavisPeabodyLeach2024RFC9562}. These strings are challenging payloads for rank transcoding because they are not natural-language messages and do not display high probability sequences. Direct subword tokenization can represent them compactly, but the resulting payload tokens may occupy low-probability positions in the model's next-token distribution\cite{Sennrich2016SubwordUnits,KudoRichardson2018SentencePiece}. A cover generator forced to select such tokens can preserve exact recoverability under matched replay while producing public text with visible artifacts or unusual statistical features. RankCloak gives an approach by which a hash-like string can be carried by generated cover tokens rather than shown as hexadecimal characters. The main methodological development of this paper is that cryptographic-artifact payloads can be treated as representations to be measured, not merely as strings to be passed through an LLM tokenizer. RankCloak therefore compares direct subword rank measurement with bounded-rank payload codecs, including ASCII-byte fixed-radix encodings and a hex-nibble encoding for lowercase hex-like payloads. This design makes the capacity to quality trade-off explicit. Direct subword ranks can be short but may impose large rank pressure, whereas bounded encodings constrain the allowed rank range while increasing the number of generated cover tokens. This representation view also separates payload recovery from public-message quality. RankCloak keeps the cover-side model, tokenizer, prompt family, deterministic rank-ordering rule, and replay procedure fixed while varying payload-side encodings and protocol variants. Non-segmented variants emit one cover token per payload rank under a single prompt context. Segmented variants split a payload into short rank chunks, emit payload-bearing forced spans across multiple public messages, and append non-payload natural tails after those forced spans. Because tails do not carry payload ranks, RankCloak reports forced-span and full-message quality proxies separately.

The resulting evaluation introduces RankCloak as a reproducible methodology for measuring how synthetic cryptographic artifacts can be represented as bounded rank sequences and expressed as LLM-generated cover text\DIFdelbegin \DIFdel{under exact-copy replay}\DIFdelend . The study \DIFdelbegin \DIFdel{is organized around how artifact-like payloads affect }\DIFdelend \DIFaddbegin \DIFadd{examines }\DIFaddend rank pressure, \DIFdelbegin \DIFdel{how bounded encodings change }\DIFdelend \DIFaddbegin \DIFadd{bounded-codec }\DIFaddend capacity and cover length, \DIFdelbegin \DIFdel{and how segmentationand natural tails change the relationship between payload-bearing spans and full messages. Exact-copy replay remains the supported recovery condition, so the sender and receiver must share the same model artifact, tokenizer, prompt rendering, filtering, deterministic configuration, codec, boundaries, and generated token identifiers. This condition is consistent with recent work showing that tokenization and retokenization mismatches can disrupt LLM-based steganography and watermarking\mbox{%DIFAUXCMD
\cite{YanMurawaki2025TokenizationInconsistency}}\hskip0pt%DIFAUXCMD
. Visible-text detokenization and retokenization is a separate, weaker recovery condition. The expanded evaluation spans public cryptographic test vectors, prompt and model families }\DIFdelend \DIFaddbegin \DIFadd{segmentation}\DIFaddend , and \DIFdelbegin \DIFdel{secondary-language conditions, and examines robustness, detectability, automated readability diagnostics, capacity, and computational overhead. RankCloak 's contribution is therefore a controlled methodology for hiding and measuring synthetic cryptographic artifacts while making representation, recovery , cover-quality, and detection trade-offs explicit. It does not establish practical robustness, security, undetectability, or human-perceived naturalness}\DIFdelend \DIFaddbegin \DIFadd{non-payload natural tails while preserving the artifact-hiding and recovery question. Saved-token exact-copy recovery is one controlled condition for exact inversion with tokenization and generation history fixed; rendered visible-text transfer and transformed-text transfer test stricter communication assumptions. The focused revision adds strict pre-feature deduplication, human-authored secondary controls, detector testing on withheld model families and with generating-model probabilities, analysis at low false-positive rates, entropy-gated embedding, and matched quantization tests. RankCloak contributes an artifact-specific hiding and recovery method with transparent empirical boundaries; it does not provide a cryptographic secrecy proof or claim universal undetectability}\DIFaddend .




\section*{Methods}\label{sec:methods}

RankCloak measures whether synthetic artifacts that do not normally look like language can be represented as LLM token ranks and then expressed as generated cover text while remaining exactly recoverable under a declared replay contract. The motivating examples are hashes, authentication tags, nonces, hexadecimal tokens, UUID identifiers, authenticated-encryption outputs, and signatures. In their original visible forms, many of these strings are recognizable by simple pattern matching (even by humans), but RankCloak asks whether the same serialized artifact can instead be carried by the token choices of an LLM-generated message. Conceptually, the method proceeds in four steps. First, the payload is converted into a sequence of integer ranks. Second, a cover prompt defines the next-token distribution used for public-text generation. Third, the sender emits the token at each required rank under that cover context. Fourth, the receiver replays the same cover context and recovers the rank of each exact saved token, following the rank-transcoding idea described by Norelli and Bronstein\cite{NorelliBronstein2026Calgacus}. If the recovered rank sequence matches the encoded payload ranks and the codec reconstructs the declared serialized bytes, the payload has been decoded exactly. The algorithms below implement two RankCloak procedures, a non-segmented procedure that maps one payload into one cover sequence and decodes the complete generated sequence, and a segmented procedure that maps a hex-like payload into several short payload-bearing spans, appends non-payload natural tails, and decodes only the forced spans. These procedures support both construction-side and detection-side measurement because they provide exact-copy recovery tests, rank-pressure measurements, cover-quality proxies, labeled cover examples, and rank-derived features that can be audited by steganalysis methods.

Sender and receiver are assumed to share a common overall LLM configuration, denoted as \(K_{\mathrm{common}}\). \DIFdelbegin \DIFdel{This configuration includes }\DIFdelend \DIFaddbegin \DIFadd{It specifies }\DIFaddend the exact model \DIFdelbegin \DIFdel{artifact}\DIFdelend \DIFaddbegin \DIFadd{revision}\DIFaddend , tokenizer, quantization and backend\DIFdelbegin \DIFdel{settings}\DIFdelend , prompt rendering, deterministic \DIFdelbegin \DIFdel{rank-ordering }\DIFdelend \DIFaddbegin \DIFadd{rank }\DIFaddend rule, payload codec, \DIFdelbegin \DIFdel{segmentation rule where used, tail and }\DIFdelend \DIFaddbegin \DIFadd{token filter, span convention, and any segmentation, tail, or }\DIFaddend lead-in \DIFdelbegin \DIFdel{policies where used, token filterwhere used, forced-span boundaries where used, and saved payload-bearing token IDs. Protocol labels select predeclared variants of this shared configuration. They are not modeled as cryptographic keys, authentication tags, signatures, or operational commands. }%DIFDELCMD <

%DIFDELCMD < %%%
\DIFdel{Recovery is evaluated under an exact-copy replay requirement for deterministic context reproduction. This means that the receiver consumes the saved generated token IDs and replays the same model, tokenizer, prompt configuration, rank-ordering rule, codec, filter, and segmentation metadata. Visible-text detokenization and retokenization is a separate, weaker replay condition. Edits, paraphrases, whitespace normalization, moderation rewrites, truncation of }\DIFdelend \DIFaddbegin \DIFadd{policy. Selecting such a configuration is analogous to selecting a common codec or protocol implementation and need not create a new per-message negotiation; a prompt phrase may be shared contextual information, but it is not treated as a formal cryptographic key. In the saved-token experimental condition, retained }\DIFaddend payload-bearing \DIFdelbegin \DIFdel{tokens, retokenization, model changes, quantization changes, prompt changes, or filter changes can alter ranks and are outside the supported exact-copy condition}\DIFdelend \DIFaddbegin \DIFadd{token IDs hold tokenization and generation history fixed while the decoder still reconstructs every contextual rank and inverts the selected payload representation; visible-text and transformed-text conditions separately test transport without preserving that original token sequence}\DIFaddend .


The exact-copy assumptions used throughout the experiments are:
\begin{itemize}
\item Shared configuration: sender and receiver use the same \DIFdelbegin \DIFdel{model artifact}\DIFdelend \DIFaddbegin \DIFadd{pinned protocol definition, including model revision}\DIFaddend , tokenizer, quantization and backend, prompt rendering, rank rule, codec, filters, \DIFaddbegin \DIFadd{span convention, }\DIFaddend segmentation policy, and tail or lead-in policy.
\item Deterministic rank ordering: ranks are one-indexed and computed from descending logits with token-ID tie-breaking.
\item Payload codecs: ASCII-byte fixed-radix coding applies to exact serialized bytes; hex-nibble coding applies only to canonical lowercase hexadecimal payloads.
\item Non-segmented decoding: the receiver replays the complete saved cover-token sequence.
\item Segmented decoding: the receiver decodes only the saved payload-bearing forced spans; natural tails are public text but carry no payload ranks.
\item Endpoint validation: rank recovery, representation recovery, serialized byte length, and SHA-256 identity are recorded separately, with exact serialized-payload equality as the primary endpoint.
\end{itemize}
The reader-facing protocol labels distinguish non-segmented variants, where recovery replays the complete generated cover-token sequence, from segmented variants, where recovery decodes only the payload-bearing forced span and ignores non-payload tails.


\subsection*{Basic rank definition}\label{sec:rank-method}

Rank-transcoding uses the rank of a token in an LLM next-token ordering as the transferable object. A rank is defined by sorting candidate tokens under a model context. Exact recovery is possible only when the sender and receiver replay the same model, tokenizer, prompt, rank rule, and generated tokens.
Let \(M\) denote an autoregressive language model with vocabulary \(\mathcal{V}\). For a token context \(h\), let \(z_M(v\mid h)\) be the next-token logit for token \(v\in\mathcal{V}\), and let \(\operatorname{id}(v)\) be its deterministic token id. The model probability is:
\begin{equation}
p_M(v\mid h)=
\frac{\exp z_M(v\mid h)}{\sum_{u\in\mathcal V}\exp z_M(u\mid h)}.
\label{eq:next-token-prob}
\end{equation}
Ranks are one-indexed; token ID breaks ties deterministically:
\begin{equation}
u \prec_h v
\Longleftrightarrow
z_M(u\mid h)>z_M(v\mid h)
\ \text{or}\
\left[z_M(u\mid h)=z_M(v\mid h)\ \text{and}\ \operatorname{id}(u)<\operatorname{id}(v)\right].
\label{eq:rank-order}
\end{equation}
The rank of a token and the inverse rank selector are,
\begin{equation}
\rho_h(v)=1+\left|\{u\in\mathcal V:u\prec_h v\}\right|,
\qquad
\operatorname{token}_h(r)=v\ \text{such that}\ \rho_h(v)=r.
\label{eq:rank-function}
\end{equation}
Here \(\operatorname{token}_h(r)\) means the token at rank \(r\) in the deterministic ordering for context \(h\). Because softmax is order-preserving, the released implementation computes this order directly from 64-bit logit values and breaks exact ties by increasing token ID.

\subsection*{Calgacus rank-transcoding baseline}\label{sec:calgacus-method}

The starting point is the Calgacus rank-transcoding construction, in which ranks rather than token identities are transferred between contexts\cite{NorelliBronstein2026Calgacus}. In the Calgacus procedure, a hidden text can be tokenized as \(e_1,e_2,\ldots\). For each hidden-text token \(e_i\), the sender computes its rank under a payload-side context \(g_i\), where \(g_i\) is determined by the preceding hidden-text tokens and any payload-side conditioning used by the protocol. The sender then emits a cover token by selecting the token with the same rank under a cover-side context \(h_i\), where \(h_i\) is determined by the cover prompt or key and the previously emitted cover tokens.

For hidden-text token \(e_i\), payload-side context \(g_i\), cover-side context \(h_i\), emitted cover token \(y_i\), and transferred rank \(r_i\), the baseline rank-transcoding step is:
\begin{equation}
r_i=\rho_{g_i}(e_i),
\qquad
y_i=\operatorname{token}_{h_i}(r_i),
\qquad
\hat r_i=\rho_{h_i}(y_i).
\label{eq:calgacus-rank-transcoding}
\end{equation}
Under exact replay with the same model, tokenizer, prompt configuration, and deterministic rank rule, the recovered rank satisfies \(\hat r_i=r_i\). The receiver can then replay the payload-side context and reconstruct the hidden-text token as \(\hat e_i=\operatorname{token}_{g_i}(\hat r_i)\). This is a same-token-length arbitrary-text baseline; one hidden-text token determines one cover token. The construction shows that a rank sequence can be carried by generated text without directly emitting the hidden tokens. Direct application to cryptographic-artifact strings, however, can impose high rank pressure. Such strings are often poorly predicted by a language model, so their payload-side tokens can occupy low-probability, high-rank positions, and selecting those same ranks under a cover prompt can degrade model-likelihood and surface diagnostics even when recovery remains exact.

RankCloak keeps the cover-side rank-transfer idea from Calgacus but changes the payload-side representation. Instead of relying only on direct LLM tokenization of the displayed artifact, RankCloak maps serialized artifact strings into explicit bounded rank sequences. In the RankCloak notation used below, the displayed artifact is \(x\), the payload codec is \(C\), and the encoded rank sequence is \(R=C(x)\). Recovery therefore proceeds by reconstructing \(\hat R\) from the cover tokens and decoding \(\hat{x}=C^{-1}(\hat R)\), rather than by replaying a natural-language hidden-text stream token by token. The direct-subword condition retains the Calgacus payload-side ranks as a separate measurement arm. The contribution relative to the baseline is the measurement framework around artifact-specific representations: direct subword rank diagnostics, bounded ASCII-byte and hex-nibble codecs, deterministic token filtering, segmented forced spans, non-payload tails, and separate forced-span versus full-message metrics.

\subsection*{RankCloak bounded payload-codec methods}\label{sec:payload-method}

RankCloak changes the payload side by applying explicit codecs to displayed synthetic payload strings. The goal is to avoid treating a hash-like or ciphertext serialization as an ordinary natural-language token sequence. Instead, the displayed artifact is converted into a controlled sequence of ranks before cover generation begins. This makes rank pressure measurable by design. A codec can trade cover length against rank magnitude so that smaller rank alphabets constrain forced token choices but require more generated tokens, while more compact representations can reduce length but may impose larger rank pressure.
For payload string \(x\), codec \(C\), encoded rank sequence \(R\), and recovered rank sequence \(\hat R\),
\begin{equation}
R=C(x)=(r_1,\ldots,r_n),
\qquad
\hat x=C^{-1}(\hat R).
\label{eq:rankcloak-codec}
\end{equation}
The direct-subword arm is not folded into Eq.~\eqref{eq:rankcloak-codec}; it tokenizes the literal serialization under the Calgacus payload context and retains the resulting ranks without imposing a fixed rank alphabet. The released implementation disables special-token insertion and accepts only a reversible tokenization whose detokenization differs from the literal UTF-8 payload by an explicitly recorded leading-space prefix. Any other affix or payload-byte change makes that representation unavailable before generation; complete framing metadata appears in Supplementary Note S1.

Non-segmented cover generation and replay recovery then use the same rank selector as the Calgacus baseline, but the ranks come from the payload codec rather than only from direct payload-side LLM token ranks:
\begin{equation}
y_i=\operatorname{token}_{h_i}(r_i),
\qquad
\hat r_i=\rho_{h_i}(y_i),
\qquad
\hat R=(\hat r_1,\ldots,\hat r_n).
\label{eq:rankcloak-nonseg}
\end{equation}
Codec-level exact recovery requires both rank recovery and an invertible codec for the displayed payload:
\begin{equation}
\hat x=x
\quad\text{when}\quad
\hat R=R
\quad\text{and}\quad
C^{-1}(C(x))=x.
\label{eq:codec-exact-recovery}
\end{equation}

Hex-nibble coding applies to canonical lowercase hexadecimal payloads. It maps each displayed hexadecimal character to one bounded rank:
\begin{equation}
\texttt{0}\mapsto 1,\quad
\texttt{1}\mapsto 2,\quad
\ldots,\quad
\texttt{9}\mapsto 10,\quad
\texttt{a}\mapsto 11,\quad
\ldots,\quad
\texttt{f}\mapsto 16.
\label{eq:hex-nibble-map}
\end{equation}
Thus each lowercase hex character maps to one rank in \(1,\ldots,16\). A 64-character SHA-256 hexadecimal serialization becomes 64 ranks under this codec. Base64 and hyphenated UUID display forms are rejected rather than silently normalized. The bounded alphabet controls maximum requested rank while increasing the number of forced positions; its relationship to observed cover diagnostics is evaluated rather than assumed.
ASCII-byte fixed-radix coding applies to arbitrary displayed payload strings. It first expresses the displayed UTF-8 byte string as base-\(B\) digits \(d_i\), then maps each digit to rank \(d_i+1\),
\begin{equation}
d_i\in\{0,\ldots,B-1\},
\qquad
r_i=d_i+1.
\label{eq:ascii-radix-map}
\end{equation}
The released codec concatenates the complete serialized byte bitstream, appends zero bits only at its end to reach a multiple of \(\log_2B\), and records the original byte length and terminal-padding count. Decoding subtracts one from each rank, validates the digit range, removes only that recorded padding, and reconstructs exactly the declared number of bytes. Consequently, the \(B=8\) condition uses \(\lceil8m/3\rceil\) ranks for \(m\) serialized bytes rather than independently padding every byte; the \(B=16\) condition uses two ranks per byte. Supplementary Note S1 gives the complete metadata and edge cases.

When a deterministic token filter is configured, ranks are computed within the allowed token set \(A_h\):
\begin{equation}
\rho^F_h(v)=1+\left|\{u\in A_h:u\prec_h v\}\right|,
\qquad
A_h\subseteq\mathcal V.
\label{eq:filtered-rank}
\end{equation}
Here \(A_h\) is the allowed token set after filtering. If no filter is configured, \(A_h=\mathcal V\). The inverse selector \(\operatorname{token}_h^F(r)\) denotes the token at rank \(r\) in the same ordered allowed set. Sender and receiver must apply the identical mask and ordering because removing one candidate shifts every lower choice, a trial is unavailable if the allowed set cannot supply a requested rank.

Algorithm~\ref{alg:nonseg} gives the non-segmented bounded-codec RankCloak procedure. The algorithm first encodes the displayed artifact into ranks. It then generates one cover token per rank under the cover prompt. Recovery replays the same prompt and exact saved token sequence, recomputes each token's rank, and decodes the recovered ranks back to the displayed payload. This algorithm is the basic measurement unit for the non-segmented cover-generation trials.

\begin{algorithm}[H]
\caption{Non-segmented RankCloak bounded-codec trial for encoding a displayed synthetic artifact as ranks, generating one cover token per rank, and recovering by exact-copy replay.}
\label{alg:nonseg}
\begin{algorithmic}[1]
\Require Synthetic payload \(x\); model \(M\); cover prompt \(p\); payload codec \(C\); optional deterministic token filter \(F\); shared configuration \(K_{\mathrm{common}}\).
\Ensure Cover text, recovered payload \(\hat{x}\), exact-recovery flag, and recorded cover metrics.
\State \(R \gets C.\Encode(x)\)
\State \(\mathrm{context} \gets \Tokenize(p)\)
\State \(\mathrm{cover\_tokens} \gets [\,]\)
\For{each rank \(r_i \in R\)}
    \State \(\mathrm{logits} \gets M.\NextLogits(\mathrm{context})\)
    \If{\(F\) is configured}
        \State \(\mathrm{allowed\_tokens} \gets F(\mathrm{logits})\)
    \Else
        \State \(\mathrm{allowed\_tokens} \gets \mathrm{vocabulary}(M)\)
    \EndIf
    \State \(y_i \gets\) token with one-indexed rank \(r_i\) among \(\mathrm{allowed\_tokens}\)
    \State append \(y_i\) to \(\mathrm{cover\_tokens}\)
    \State append \(y_i\) to \(\mathrm{context}\)
\EndFor
\State \(\mathrm{cover\_text} \gets \Detokenize(\mathrm{cover\_tokens})\)
\Statex
\State \(\mathrm{context} \gets \Tokenize(p)\)
\State \(\mathrm{recovered\_ranks} \gets [\,]\)
\For{each exact saved token \(y_i\) in \(\mathrm{cover\_tokens}\)}
    \State \(\mathrm{logits} \gets M.\NextLogits(\mathrm{context})\)
    \If{\(F\) is configured}
        \State \(\mathrm{allowed\_tokens} \gets F(\mathrm{logits})\)
    \Else
        \State \(\mathrm{allowed\_tokens} \gets \mathrm{vocabulary}(M)\)
    \EndIf
    \State append \(\RankOf(y_i,\mathrm{logits},\mathrm{allowed\_tokens})\) to \(\mathrm{recovered\_ranks}\)
    \State append \(y_i\) to \(\mathrm{context}\)
\EndFor
\State \(\hat{x} \gets C.\Decode(\mathrm{recovered\_ranks})\)
\State \(\mathrm{exact\_recovery} \gets (\hat{x}=x)\)
\State \Return \(\mathrm{cover\_text}, \hat{x}, \mathrm{exact\_recovery}\)
\end{algorithmic}
\end{algorithm}

Before Algorithm~\ref{alg:nonseg} is executed, the selected representation validates payload syntax, exact byte length, codec applicability, and the source digest. The retained record saves the requested ranks, cover-token IDs, prompt and configuration identities, codec metadata, and failure position where applicable. After decoding, exact rank and representation equality are distinguished from exact serialized-payload equality, which additionally requires the declared byte length, bytes, syntax, and SHA-256 digest to match. The same generation and replay loop also applies when the validated direct-subword representation supplies \(R\); its payload-side reconstruction and framing checks are detailed in Supplementary Note S1.

\subsection*{RankCloak segmented multi-cover method}\label{sec:encoding-method}

The segmented multi-cover method extends RankCloak from one cover sequence to several shorter public messages. Long sequences of forced ranks can accumulate topic drift, formatting artifacts, or low-probability continuations. Segmentation restarts the cover context for each chunk, allowing the effect of distributing the payload across several messages to be measured without claiming that those messages are human-natural. The method is also useful analytically because it separates the payload-bearing forced span from the rest of the public message.

For a full rank sequence \(R\), segmentation splits the payload into \(m\) chunks:
\begin{equation}
R = R^{(1)}\Vert R^{(2)}\Vert \cdots \Vert R^{(m)}.
\label{eq:rank-segmentation}
\end{equation}
Here \(\Vert\) denotes sequence concatenation, and each \(R^{(j)}\) is an ordered, non-overlapping rank chunk. The completed design includes frozen segmented schedules and an ablation over 4-, 8-, 16-, and 32-rank chunks rather than treating one chunk size as universally preferred.

For segment \(j\), the forced span emits one token per payload rank. Recovery uses only those forced tokens:
\begin{equation}
y^{(j)}_k=
\operatorname{token}_{h^{(j)}_k}\left(r^{(j)}_k\right),
\qquad
\hat r^{(j)}_k=
\rho_{h^{(j)}_k}\left(y^{(j)}_k\right).
\label{eq:segmented-forced-span}
\end{equation}
When a filter is configured, both operations use the filtered ordering defined after Eq.~\eqref{eq:filtered-rank}. After the forced span \(y^{(j)}\), the sender appends a greedy non-payload tail \(z^{(j)}\). The public message is the tokenizer rendering of the concatenated token sequence \(y^{(j)}\Vert z^{(j)}\). The receiver ignores \(z^{(j)}\) and decodes only \(y^{(j)}\):
\begin{equation}
m_j=\operatorname{text}\left(y^{(j)}\Vert z^{(j)}\right),
\qquad
\hat R=\hat R^{(1)}\Vert\cdots\Vert\hat R^{(m)}.
\label{eq:segmented-message}
\end{equation}
Here \(\operatorname{text}(\cdot)\) means ordinary tokenizer rendering of the token sequence; it is not an additional cryptographic operation. The canonical expression in Eq.~\eqref{eq:segmented-message} has no lead-in. For the evaluated lead-in extension, a non-payload prefix is generated before \(y^{(j)}\), and the record stores exact lead-in, forced, tail, and full token-ID arrays, a token-role mask, and half-open forced-span offsets. The decoder uses those declared offsets and does not infer a payload boundary from visible prose.

A core component of this method is the forced-span and tail separation. The payload is carried only by the forced span, while the tail is a greedy continuation whose effect on the complete public message is measured separately. This distinction prevents a misleading quality result: full-message diagnostics can change because of the non-payload tail even if the payload-bearing forced span remains low probability or exhibits surface artifacts. For that reason, the Results report forced-span metrics and full-message metrics separately.

Algorithm~\ref{alg:segmented} describes the segmented multi-cover protocol. It encodes the hex-like payload, splits the rank sequence into short chunks, generates one forced span per chunk, appends the selected tail-policy output after each forced span, and recovers the payload by replaying only the forced spans. Each segment is its own replay component: generation starts from its declared prompt, and recovery reads the exact saved token IDs within its declared forced-span boundary.

\begin{algorithm}[H]
\caption{Segmented multi-cover RankCloak trial for distributing a hex-nibble payload across payload-bearing forced spans with non-payload natural tails.}
\label{alg:segmented}
\begin{algorithmic}[1]
\Require Synthetic hex-like payload \(x\); model \(M\); hex-nibble codec \(C_{\mathrm{hex}}\); segment size \(s\); prompt schedule \(P=(p_1,\ldots,p_m)\); deterministic token filter \(F\); tail policy \(T\); shared configuration \(K_{\mathrm{common}}\).
\Ensure Public messages \((\mathrm{message}_1,\ldots,\mathrm{message}_m)\), recovered payload \(\hat{x}\), exact-recovery flag, forced-span metrics, and full-message metrics.
\State \(R \gets C_{\mathrm{hex}}.\Encode(x)\)
\State split \(R\) into chunks \((R_1,\ldots,R_m)\), each with at most \(s\) ranks
\For{each segment \(j=1,\ldots,m\)}
    \State \(\mathrm{context} \gets \Tokenize(P[j])\)
    \State \(\mathrm{forced\_tokens}_j \gets [\,]\)
    \For{each rank \(r \in R_j\)}
        \State \(\mathrm{logits} \gets M.\NextLogits(\mathrm{context})\)
        \State \(\mathrm{allowed\_tokens} \gets F(\mathrm{logits})\)
        \State \(y \gets\) token with one-indexed rank \(r\) among \(\mathrm{allowed\_tokens}\)
        \State append \(y\) to \(\mathrm{forced\_tokens}_j\)
        \State append \(y\) to \(\mathrm{context}\)
    \EndFor
    \State \(\mathrm{tail\_tokens}_j \gets \Call{TailPolicy}{M,\mathrm{context},F,T}\)
    \State \(\mathrm{message}_j \gets \Detokenize(\operatorname{concat}(\mathrm{forced\_tokens}_j,\mathrm{tail\_tokens}_j))\)
\EndFor
\State emit public messages \((\mathrm{message}_1,\ldots,\mathrm{message}_m)\)
\Statex
\State \(\mathrm{recovered\_chunks} \gets [\,]\)
\For{each segment \(j=1,\ldots,m\)}
    \State \(\mathrm{context} \gets \Tokenize(P[j])\)
    \State read the exact saved tokens within the declared forced-span boundary for \(R_j\)
    \State ignore all tail tokens
    \State \(\mathrm{segment\_ranks} \gets [\,]\)
    \For{each forced token \(y\) in the payload-bearing span}
        \State \(\mathrm{logits} \gets M.\NextLogits(\mathrm{context})\)
        \State \(\mathrm{allowed\_tokens} \gets F(\mathrm{logits})\)
        \State append \(\RankOf(y,\mathrm{logits},\mathrm{allowed\_tokens})\) to \(\mathrm{segment\_ranks}\)
        \State append \(y\) to \(\mathrm{context}\)
    \EndFor
    \State append \(\mathrm{segment\_ranks}\) to \(\mathrm{recovered\_chunks}\)
\EndFor
\State \(\hat{x} \gets C_{\mathrm{hex}}.\Decode(\mathrm{concat}(\mathrm{recovered\_chunks}))\)
\State \(\mathrm{exact\_recovery} \gets (\hat{x}=x)\)
\State \Return \((\mathrm{message}_1,\ldots,\mathrm{message}_m),\hat{x},\mathrm{exact\_recovery}\)
\end{algorithmic}
\end{algorithm}

The non-lead-in form in Algorithm~\ref{alg:segmented} supplies the canonical segment layout. The lead-in extension inserts a rank-1 prefix before the forced span; its exact token IDs must be replayed before forced-rank recovery, and its saved length shifts the half-open forced boundary. Lead-in lengths 2, 4, 8, 16, and 32 were compared with the no-lead-in condition. Single-topic segmentation reuses its assigned prompt, whereas multi-topic segmentation rotates deterministically through the declared prompt categories. Segment order, prompts, boundaries, and token roles are retained metadata rather than information inferred by the decoder.

Algorithm~\ref{alg:tail} implements the fixed natural-tail policy available to Algorithm~\ref{alg:segmented}. The tail is generated greedily, begins only after the payload-bearing forced span, and is not decoded. Its purpose is to measure how a non-payload continuation changes the complete public message while preserving a separate measurement of the forced span. This tail extension is a sentence-completion step rather than part of the hidden payload. A forced span can end abruptly because its tokens were selected to carry payload ranks, not necessarily to finish a phrase or sentence. The minimum length and sentence-boundary rule allow a continuation before the hard cap prevents non-payload text from growing without bound; no claim of human-perceived completeness follows from this proposed heuristic.

\begin{algorithm}[H]
\caption{Greedy natural-tail policy used after each segmented forced span; the generated tail carries no payload ranks and is ignored by the decoder.}
\label{alg:tail}
\begin{algorithmic}[1]
\Require Model \(M\); current context after the forced span; deterministic token filter \(F\).
\Ensure Non-payload tail tokens.
\State \(\mathrm{tail\_tokens} \gets [\,]\)
\While{\(|\mathrm{tail\_tokens}| < 60\)}
    \State choose the rank-1 token \(z\) under \(M.\NextLogits(\mathrm{context})\) among tokens allowed by \(F\)
    \State append \(z\) to \(\mathrm{tail\_tokens}\)
    \State append \(z\) to \(\mathrm{context}\)
    \If{\(|\mathrm{tail\_tokens}| \geq 20\) and \(\Detokenize(\mathrm{tail\_tokens})\) ends at a sentence boundary}
        \State \textbf{break}
    \EndIf
\EndWhile
\State \Return \(\mathrm{tail\_tokens}\)
\end{algorithmic}
\end{algorithm}

Algorithm~\ref{alg:tail} realizes \textsc{TailPolicy} as a fixed sentence completion of \(20\) to \(60\) tokens. The completed primary matrix uses a separately frozen dynamic condition as its reference level: after at least eight greedy tokens, stopping requires terminal punctuation or a double newline, balanced delimiters, and no declared dangling connective or punctuation, with a 256-token emergency cap recorded as censoring. A no-tail condition returns an empty tail and all three policies are non-payload conditions ignored during decoding (their complete stopping definitions appear in Supplementary Note S1).

\subsection*{Evaluation corpus, controls, and metrics}\label{sec:design-method}\label{sec:quality-method}

The experiments use deterministic synthetic payloads that resemble cryptographic or operational artifacts while avoiding real secret material. The corpus contains 480 public deterministic test vectors, 60 in each of eight algorithm-defined classes; SHA-256 hexadecimal digests, HMAC-SHA256 hexadecimal tags, deterministic 96-bit nonces, 128-bit hexadecimal tokens, UUIDv4 identifiers, AES-256-GCM base64 outputs, ChaCha20-Poly1305 base64 outputs, and Ed25519 base64 signatures\cite{NIST2015SecureHashStandard,Krawczyk1997RFC2104,NIST2007GCM,NirLangley2018RFC8439,JosefssonLiusvaara2017RFC8032,Josefsson2006RFC4648,DavisPeabodyLeach2024RFC9562}. Authenticated-encryption and signature rows are genuine outputs of the declared algorithms under reproducible public test inputs. No real credentials, API keys, private keys, account identifiers, operational secrets, commands, private user data, or deployed traffic are used.

Each payload is evaluated first as a representation problem and then, where applicable, as a cover-generation problem. The representation diagnostics measure how the exact serialization behaves under direct subword tokenization and how many ranks are required by the bounded codecs. Cover-generation conditions include direct subword ranks, ASCII-byte fixed-radix coding with \(B=8\), ASCII-byte fixed-radix coding with \(B=16\), raw hex-nibble coding for eligible lowercase hexadecimal payloads, and single- and multi-topic segmented hex-nibble protocols. All generation, recovery, and diagnostic rows use the deterministic rank-ordering rule in Eqs.~\eqref{eq:rank-order} and \eqref{eq:rank-function} as part of the shared replay configuration.

Generation used Meta-Llama-3-8B-Instruct, Qwen2.5-7B-Instruct, and Mistral-7B-Instruct-v0.3 in pinned Q4\_K\_M configurations\DIFaddbegin \DIFadd{, an approximately 4-bit quantized format, }\DIFaddend with their embedded GGUF tokenizers. Eighteen English templates span six prompt categories. Secondary Spanish and Simplified Chinese conditions evaluate saved-token recovery but do not establish multilingual human naturalness. The primary matrix contains 14,400 generated units; 6,480 RankCloak payload trials and 7,920 prompt, model, representation, and length-matched ordinary controls. The complete ledger reconciles 38,504 declared work units across the primary, robustness, ablation, multilingual, evaluator, and detector stages. Complete corpus construction, model and tokenizer identities and hashes, prompt schedules, and work-unit accounting appear in Supplementary Note S2.

The protocol variants define how encoded ranks become public text. Non-segmented trials use Algorithm~\ref{alg:nonseg} and replay the complete saved cover-token sequence. Segmented trials use Algorithm~\ref{alg:segmented}; Algorithm~\ref{alg:tail} or another declared tail condition follows each forced span, while recovery uses only the saved forced tokens. Prompt-matched controls provide comparison text for automated diagnostics and detector training; they are not human-written naturalness baselines. The deterministic safe-text filter rejects declared control, replacement, code, markup, URL, escape, and formatting fragments. The ablation compares it with no filter and an isolated roundtrip-stable filter, lead-in lengths 0, 2, 4, 8, 16, and 32, chunk sizes 4, 8, 16, and 32, dynamic, fixed, and no-tail policies, and frozen single- versus multi-topic schedules.

The Results report exact serialized-payload recovery, representation and rank diagnostics, cover length, requested-rank pressure, token log probability, deterministic surface artifacts, and automated readability diagnostics. A balanced 504-stimulus subset covers seven generation conditions and all 18 English templates; its surface measures include Flesch Reading Ease\cite{Flesch1948Readability}, frozen artifact flags, and TF-IDF prompt similarity, with intervals bootstrapped over prompt templates. Source-model and held-out-evaluator likelihoods are reported as model diagnostics, not participant ratings. No human study was performed.

For segmented rows, every cover-quality measurement is assigned to one of two scopes. Forced-span metrics are computed only over payload-bearing tokens decoded by the receiver. Full-message metrics are computed over the complete rendered message, including non-payload tails and any experimental lead-in. This scope distinction is necessary because tails can change full-message proxies while carrying no payload ranks. Declared unavailable combinations are excluded from the applicable estimand and reported separately rather than classified as observed failures. Complete metric definitions, transformation coverage, ablation schedules, and exclusion rules appear in Supplementary Notes S3 through S6.

For paired segmented trials, tail gain is the full-message mean token log probability minus the forced-span mean token log probability, \(G_{\mathrm{tail}}=\bar\ell_{\mathrm{full}}-\bar\ell_{\mathrm{forced}}\). Positive values mean that averaging in the non-payload continuation raises the full-message model score; they do not imply a change in the payload-bearing span. The empirical frontier reports forced-span and full-message views separately on the four-class common hexadecimal support, with payload-bootstrap uncertainty and automated surface-flag burden detailed in Supplementary Note S5.

\subsection*{Recovery, robustness, and capacity endpoints}\label{sec:replay-method}\label{sec:capacity-method}\label{sec:endpoint-method}

The \DIFdelbegin \DIFdel{strongest recovery endpoint passes the original saved token IDs }\DIFdelend \DIFaddbegin \DIFadd{saved-token endpoint supplies the retained payload-bearing token sequence }\DIFaddend to the decoder under \(K_{\mathrm{common}}\)\DIFdelbegin \DIFdel{. Saved IDs are an exact-copy record, not an ordinary text channel}\DIFdelend \DIFaddbegin \DIFadd{, thereby testing exact contextual rank reconstruction and codec inversion while tokenization fidelity and generation history are held fixed; retaining IDs does not make those inverse computations unnecessary or expose the payload serialization}\DIFaddend . Visible-text replay \DIFdelbegin \DIFdel{first }\DIFdelend \DIFaddbegin \DIFadd{instead }\DIFaddend detokenizes and retokenizes the rendered string, \DIFdelbegin \DIFdel{so an apparently unchanged string can yield different IDs or boundaries\mbox{%DIFAUXCMD
\cite{YanMurawaki2025TokenizationInconsistency}}\hskip0pt%DIFAUXCMD
. Arbitrary editing, paraphrasing, formatting changes, truncation of payload-bearing tokens, cross-model decoding, and prompt or filter search are not supported recovery modes.
Controlled secondary tests separately evaluate unmodified visible text, }\DIFdelend \DIFaddbegin \DIFadd{and controlled transformed-text tests additionally apply }\DIFaddend final-tail truncation, Unicode NFKC, quote conversion, whitespace and line-ending changes, deterministic character and token edits, markdown copy/paste, paraphrase, limited canonicalization, \DIFdelbegin \DIFdel{and }\DIFdelend \DIFaddbegin \DIFadd{or }\DIFaddend cross-model decoding\DIFdelbegin \DIFdel{. Transformations are frozen as separate conditionsrather than composed into an uncontrolled channel}\DIFdelend \DIFaddbegin \DIFadd{\mbox{%DIFAUXCMD
\cite{YanMurawaki2025TokenizationInconsistency}}\hskip0pt%DIFAUXCMD
. These are distinct transmission conditions: token-preserving delivery can retain the exact token stream, ordinary rendered-text delivery does not, and paraphrased delivery changes the visible string}\DIFaddend . Final-tail truncation removes only declared non-payload tokens and is not evidence of arbitrary truncation tolerance\DIFdelbegin \DIFdel{. First-divergence }\DIFdelend \DIFaddbegin \DIFadd{; first-divergence }\DIFaddend labels localize the first recorded mismatch but do not establish its cause.


For a bounded representation containing \(H\) source bits and a rank alphabet of size \(B\), the minimum fixed-radix forced length and nominal forced-position rate are
\begin{equation}
n_B=\left\lceil\frac{H}{\log_2B}\right\rceil,
\qquad
R_B=\frac{H}{n_B}\le\log_2B.
\label{eq:capacity}
\end{equation}
Here \(H\) is the information actually presented to that codec, eight bits per exact serialized byte for the byte codecs and four bits per canonical hexadecimal character for the nibble codec. For the byte bitstream, \(n_B=\lceil8m/\log_2B\rceil\), consistent with Eq.~\eqref{eq:ascii-radix-map}. Direct subword ranks do not have a fixed \(B\) and are excluded from this bounded-capacity identity.

For representation comparisons on a common artifact support, we additionally report effective artifact bits per forced token, \(R_{\mathrm{art}}=H_{\mathrm{art}}/n_{\mathrm{forced}}\). Here \(H_{\mathrm{art}}\) is held fixed across representations, four bits per canonical hexadecimal character in the common-hex analysis, rather than set to the codec's serialized input width. This descriptive rate permits direct and bounded representations to be compared on the same artifacts; it is not a cryptographic-security capacity.

If segment \(j\) has \(\ell_j\) lead-in tokens and \(t_j\) tail tokens, and \(\delta\ge0\) denotes any other counted non-payload extension under the retained accounting convention, full cover length and effective rate are,
\begin{equation}
N_{\mathrm{full}}
=n_{\mathrm{forced}}+\sum_{j=1}^{J}(\ell_j+t_j)+\delta,
\qquad
R_{\mathrm{eff}}=\frac{H}{N_{\mathrm{full}}}
\le\frac{H}{n_{\mathrm{forced}}}=R_B
\quad(n_{\mathrm{forced}}=n_B).
\label{eq:effective-rate}
\end{equation}
Segmentation changes context resets and auxiliary text, not the number of forced codec positions. Positive cover-length residuals are overhead rather than payload capacity.

For eligible-token probabilities \(p_{(1)}(h)\ge\cdots\ge p_{(B)}(h)\), any bounded selection \(r\le B\) obeys the same-context conditional quality relation
\begin{equation}
-\log p_{(1)}(h)
\le-\log p_{(r)}(h)
\le-\log p_{(B)}(h),\qquad
Q_B=\frac{1}{n_B}\sum_{i=1}^{n_B}-\log p_M(y_i\mid h_i).
\label{eq:quality-bound}
\end{equation}
Some tables report mean log probability \(\bar\ell=-Q_B\), for which larger values are less negative. These bounds concern model likelihood under identical filtered histories; they are not statements about grammar, meaning, human naturalness, security, or detectability.

The confirmatory endpoint is exact equality between the decoded serialized payload bytes and the declared source, including length, syntax, and SHA-256 digest. Exact rank and representation equality are retained as diagnostics. Trial-level Wilson 95\% intervals are reported\cite{Wilson1927Interval}; a strict payload group passes only if every constituent prompt, model, and eligible protocol trial succeeds, and a group defined by model and payload applies the same rule within one model. Segments are never counted as independent payload trials. Grouped bootstrap procedures use payload or prompt template as the resampling unit according to the estimand. Declared unavailable combinations never enter a failure denominator.

\subsection*{Neural steganalysis, statistical analysis, and computational overhead}\label{sec:detector-method}\label{sec:statistics-method}

\DIFdelbegin \DIFdel{The frozen detector corpus contains }\DIFdelend \DIFaddbegin \DIFadd{Strict corpus preparation preceded feature extraction and tokenization. Unicode NFKC normalization, case folding, Unicode-whitespace collapse, and hashing identified one normalized exact-duplicate group; its complete matched pair was removed from the }\DIFaddend 15,\DIFdelbegin \DIFdel{840 rows, balanced between 7, 920 RankCloak covers and 7}\DIFdelend \DIFaddbegin \DIFadd{840-row source corpus, leaving 15,838 observations. Declared near duplicates had cosine similarity at least 0.95 under L2-normalized character-boundary TF-IDF with 3--5 character \(n\)-grams and minimum document frequency two. Payload groups connected by an exact or near-duplicate edge formed immutable components assigned to 9,904 training, 2}\DIFaddend ,\DIFdelbegin \DIFdel{920 matched ordinary controls and grouped by the 480 payloads. A TextCNN-equivalent classifier\mbox{%DIFAUXCMD
\cite{Kim2014CNN,Wen2019CNNTextSteganalysis} }\hskip0pt%DIFAUXCMD
uses lowercase tokens, a 30}\DIFdelend \DIFaddbegin \DIFadd{952 validation, and 2}\DIFaddend ,\DIFdelbegin \DIFdel{000-term vocabulary,length 256, 300-dimensional embeddings,100 filters at widths 3 and 5, dropout 0.5, and Adam learning rate 0.001. A DeBERTa-v3-base classifier uses length 256,learning rate \(2\times10^{-5}\), weight decay 0.01, and three epochs\mbox{%DIFAUXCMD
\cite{He2021DeBERTa}}\hskip0pt%DIFAUXCMD
. Payload groups, never individual rows, define train/test splits}\DIFdelend \DIFaddbegin \DIFadd{982 test observations. The audit found zero payload, pair, normalized-text, component, or declared near-duplicate links across partitions; it cannot identify semantic paraphrases below the normalization and threshold rule}\DIFaddend .

\DIFdelbegin \DIFdel{Each architecture is fit on one matched split, 18 held-out-template splits, three leave-one-model splits, and six leave-one-codec splits, for 56 completed fits and 55, 440 predictions.
The split definitions and seeds }\DIFdelend \DIFaddbegin \DIFadd{TextCNN and a pinned DeBERTa-v3-base revision were trained specifically on RankCloak versus matched ordinary-model controls using the predeclared architectures. An exact-model-aware logistic detector, representing an attacker with likelihood information from the same model that generated the text, used summary surprisal features computed from saved token log probabilities for both labels, with regularization selected by validation ROC-AUC; ordinary-control ranks were unavailable and therefore excluded. Leave-one-family-out fits excluded the target model family from both training and validation. A secondary test selected 250 human-written responses from the Databricks Dolly corpus across 16 available templates without using human labels for fitting or threshold selection; it is a corpus comparison, not a participant study or naturalness rating.
}

\DIFadd{ROC-AUC used tie-aware empirical ranks. Thresholds for low false-positive-rate operating points (low-FPR) maximized validation sensitivity subject to the exact allowed false-positive count, were conservatively tie-broken, and }\DIFaddend were frozen before \DIFdelbegin \DIFdel{fitting. Primary endpoints are ROC-AUC and balanced accuracy; precision,Brier score, and sensitivity at }\DIFdelend \DIFaddbegin \DIFadd{test evaluation; empirical }\DIFaddend 1\DIFdelbegin \DIFdel{\% false-positive rate are exploratory. Payload-group bootstrap intervals use }\DIFdelend \DIFaddbegin \DIFadd{\% or 0.1\% results required at least 100 or 1,000 validation and test negatives, respectively, and unavailable estimates were not interpolated. Confidence intervals used }\DIFaddend 2,000 resamples \DIFdelbegin \DIFdel{\mbox{%DIFAUXCMD
\cite{Efron1979Bootstrap}}\hskip0pt%DIFAUXCMD
. Audits test exact visible-text }\DIFdelend \DIFaddbegin \DIFadd{of complete deduplication clusters. The entropy extension fixed ungated, median-gate, and 75th-percentile strict-gate thresholds from ordinary top-\(p\) development traces before evaluation and consumed payload ranks only at eligible positions. Its fixed-payload capacity is conditional on completion. The paired quantization experiment held one pinned Qwen2.5-7B-Instruct revision, tokenizer, prompt, payload, codec, filter, target length, sampler, and backend fixed while comparing Q4\_K\_M and Q8\_0, approximately 4-bit }\DIFaddend and \DIFdelbegin \DIFdel{token-sequence identity, then diagnose normalized character \(n\)-gram near-duplicates at the declared threshold. Excluding affected test payload groups without retraining is a sensitivity check and does not eliminate the near-duplicate limitation. Fixed-seed repeats establish same-device CUDA repeatability only, CPU/GPU equivalence was not tested}\DIFdelend \DIFaddbegin \DIFadd{8-bit quantized versions of the same pinned Qwen model; paired recovery, grouped bootstrap, exact McNemar, and bidirectional cross-quantization detector transfer were prespecified. Supplementary Notes S11--S15 give the complete calibration, partition, threshold, count, and provenance details}\DIFaddend .


Statistical inference respects payload and prompt-template grouping, with segments nested within covers. Paired ablations use payload as the pairing and bootstrap unit. Mixed models include payload and prompt grouping where supported, negative-binomial models for overdispersed artifact counts, and Gaussian models for continuous outcomes\cite{Bates2015lme4,Brooks2017glmmTMB}. Holm adjustment controls selected contrast families; Benjamini-Hochberg values accompany exploratory families. Convergence, dispersion, zero-inflation, singularity, and residual diagnostics are retained. Three of five fitted mixed models were singular, no undeclared fixed-effects fallback was substituted, and compact ablation, topic, auxiliary-detector, and near-duplicate analyses remain labeled exploratory or post-outcome where applicable. Complete model formulas, bootstrap seeds, split definitions, diagnostics, and multiplicity details are in Supplementary Notes S5 through S7.

Generation overhead uses inclusive wrapper wall time retained by the harness, including the reported setup, generation, and supported replay scopes. Payload throughput is representation bits divided by generation time, and effective rate follows Eq.~\eqref{eq:effective-rate}. Detector benchmarks report fixed-device CUDA wall/training time and peak allocated VRAM. CPU time, repeated warm-up measurements, and cross-device hardware generalization were unavailable and were not inferred. Complete timing cells and component scopes appear in Supplementary Note S9.




\section*{Results}\label{sec:results}

\subsection*{Expanded evaluation and exact-copy recovery}\label{sec:results-recovery}

The completed primary design produced 14,400 work units across three model families and 18 English prompt templates. Of these 6,480 were payload bearing RankCloak trials spanning the six protocol/codec configurations and 7,920 were ordinary controls. All 6,480 decoded serialized payloads exactly under saved-token replay (1.000000; Wilson 95\% CI \([0.999408, 1.000000]\); Table~\ref{tab:design-recovery}). The stricter aggregation also passed for all 480 payload groups (CI \([0.992061, 1.000000]\)) and all 1,440 groups defined by model and payload (CI \([0.997339, 1.000000]\)). These grouped endpoints prevent repeated prompts and protocols from inflating the apparent number of independently generated artifacts.

The recovery question was whether codec inversion remained exact after forcing the ranks into cover contexts, rather than whether emitted ranks could merely be read back from the same loop. Success was defined at the original serialization and digest, and every model and eligible protocol stratum reached that endpoint. The all-success result therefore supports deterministic implementation fidelity across the evaluated prompt/model matrix. Its uncertainty is one-sided in substance; the data bound the rate of error under this completed design but cannot show that an unobserved replay environment or transmission channel would also be error-free.

Secondary exact-copy tests recovered 288/288 Spanish trials and 288/288 Simplified Chinese trials, or 576/576 in total. They establish recovery for those retained prompts under the same saved-token contract; they do not establish multilingual naturalness or robustness to ordinary text transfer. Across primary, ablation, robustness, multilingual, evaluator, and detector work, the final ledger reconciled 38,504/38,504 units: 38,072 succeeded and 432 were declared unavailable. No unit remained and none was classified as an execution failure. Recovery strata and secondary-language endpoints appear in Supplementary Tables S5 and S6.

\begin{table}[H]
\centering
\small
\begin{tabularx}{\textwidth}{L{0.25\textwidth}Y L{0.24\textwidth}}
\toprule
Design or endpoint & Scope & Result \\
\midrule
Payload corpus & Eight algorithm-defined classes, 60 public deterministic artifacts per class & 480 artifacts \\
Model and prompt coverage & Three open-source 7B to 8B families; 18 English templates in six categories & 14,400 primary work units (6,480 payload; 7,920 control) \\
Primary exact-copy recovery & Serialized payload, saved-token replay & 6,480/6,480; 1.000000 (95\% CI \([0.999408, 1.000000]\)) \\
Strict grouped sensitivity & Payload groups; groups defined by model and payload & 480/480 (CI \([0.992061, 1.000000]\)); 1,440/1,440 (CI \([0.997339, 1.000000]\)) \\
Spanish secondary recovery & Saved-token replay & 288/288 (CI \([0.986837, 1.000000]\)) \\
Simplified Chinese secondary recovery & Saved-token replay & 288/288 (CI \([0.986837, 1.000000]\)) \\
Complete work ledger & All declared revision matrices & 38,504/38,504 reconciled: 38,072 success; 432 unavailable; 0 execution failure \\
\bottomrule
\end{tabularx}
\caption{\label{tab:design-recovery}\textbf{Evaluation design and recovery endpoints.} Confidence intervals are trial- or group-level Wilson 95\% intervals as identified. Recovery requires saved-token exact-copy replay with the declared shared configuration. Unavailable work units are configuration/retained-output absences, not observed decoding failures.}
\end{table}

\subsection*{Payload representation, rank pressure, and capacity tradeoff}\label{sec:results-capacity}

Figure~\ref{fig:capacity} compares payload representations across 480 unique artifacts in eight classes and three models. The 1,440 direct observations across payloads and models contained 60,861 direct-rank positions. Direct-subword representations required 11 to 77 forced symbols, with class medians ranging from 20 to 66, and the largest observed direct rank was 145,498. These counts and ranks depend on the model and tokenizer rather than defining a fixed alphabet.
The deterministic codecs instead imposed rank ceilings of 8 or 16 while generally increasing forced length. On the four-class common hexadecimal support, median effective artifact rates were approximately 4.571 for direct subword, 4.000 for hex nibble, 2.000 for ASCII \(B=16\), and 1.497 for ASCII \(B=8\) artifact bits per forced token. These descriptive rates compare representations on identical eligible artifacts; they are not cryptographic-security capacities.
The whiskers in Fig.~\ref{fig:capacity} are observed 5th to 95th percentiles, not confidence intervals. The detailed theoretical capacity and tail validation tests a different question: whether retained records satisfy the declared algebraic identities and conditional bounds (in Supplementary Fig. S6 and Table S14). The common-support empirical capacity and quality frontier follows in Fig.~\ref{fig:empirical-frontier}.

\begin{figure}[H]
\centering
\includegraphics[width=0.97\textwidth]{figures/payload_representation_rank_capacity.pdf}
\caption{\label{fig:capacity}\textbf{Payload representation, rank pressure, and capacity tradeoff.} Direct-subword symbol counts and rank pressure are summarized for 1,440 observations across payloads and models comprising 60,861 direct-rank positions. Their values are tokenizer and model dependent. Bounded codecs constrain requested ranks to at most 8 or 16, at the cost of more forced positions. Effective artifact-rate comparisons use the four-class common hexadecimal support and are descriptive rather than cryptographic-security capacities. Whiskers show observed 5th to 95th percentiles, not confidence intervals.}
\end{figure}

Figure~\ref{fig:empirical-frontier} compares mean token count with mean source-model token log probability over 4,320 trials on the four-class common hexadecimal support. Three models and six protocols form 18 model and protocol cells, and separate forced-span and full-message views give 36 plotted cells. Compact direct-subword spans generally occupy the low-length region but can have poorer source-model likelihood, whereas bounded protocols trade longer forced spans for controlled rank alphabets. Segmented full-message points move toward better full-message likelihood by adding non-payload natural-tail length; this movement is not improved likelihood on the payload-bearing forced span. Higher mean token log probability is better; point area represents mean automated surface-flag burden, and uncertainty uses 2,000-resample payload-bootstrap 95\% confidence intervals. The complete primary matrix contains no lead-in condition, so exploratory lead-in rows were not pooled. These automated diagnostics are not human-perceived judgments, and the empirical frontier is distinct from the theoretical capacity and tail validation.

\begin{figure}[H]
\centering
\includegraphics[width=0.97\textwidth]{figures/empirical_capacity_quality_frontier.pdf}
\caption{\label{fig:empirical-frontier}\textbf{Empirical capacity and quality frontier.} Mean token count is shown against mean source-model token log probability for 4,320 trials on four-class common hexadecimal support, spanning three models, six protocols, 18 model and protocol cells, and 36 forced-span or full-message cells. Higher mean token log probability is better; point area represents mean automated surface-flag burden. Error bars are 2,000-resample payload-bootstrap 95\% confidence intervals. The primary matrix contains no lead-in condition. Likelihood and surface flags are automated diagnostics, not human judgments, and this empirical frontier is distinct from the theoretical capacity and tail validation in Supplementary Fig. S6.}
\end{figure}

\subsection*{Exact-copy replay and transmission fragility}\label{sec:results-robustness}

Replay mode materially changed recovery (Fig.~\ref{fig:robustness}). Saved-token replay recovered 144/144 robustness covers, whereas detokenizing the same covers and retokenizing their visible text recovered 88/144 (0.611111). So even an unedited visible string did not guarantee the original token sequence. Raw unmodified transmission recovered 144/144, but common transformations were damaging. Unicode NFKC recovered 82/144 (0.569444), quote conversion 56/144 (0.388889), and whitespace trimming 28/144 (0.194444). Markdown copy/paste and paraphrase recovered 0/144, and the two cross-model decoder outcomes for each cover recovered 0/288.

These conditions address a different question from the primary result of whether a cover can leave its saved-token record and still recreate the forced ranks. The answer depended sharply on the frozen channel definition. The saved-ID, visible-retokenization, and raw-transmission rows are distinct replay procedures and should not be pooled. Their contrast shows why the receiver contract must name the tokenization path, not merely say that the visible text was copied exactly. Cross-model zero recovery further shows that similar prompt intent does not substitute for model and tokenizer identity.

Deleting only the declared final 10\% of a message recovered 144/144 because that operation was constrained to non-payload tail tokens. It is not evidence of arbitrary truncation tolerance; removing a forced token or earlier context would change the decoded sequence. The broader matrix also showed low recovery after token or character edits and line-ending conversion (Supplementary Note S4). Limited canonicalization preserved unmodified controls but did not recover the available transformed cells, so it did not demonstrate a general mitigation.

Tail-only truncation is consequently a boundary check on the decoder definition rather than a successful error-correction result. The decoder is designed to ignore tails, so removing only tail labels should be harmless. By contrast, prefix changes and edits inside a forced span can alter a segment's initial tokenization or model history and propagate. The failed limited-canonicalization cells show that normalizing a small declared set after the fact was insufficient in this matrix; no broader mitigation is inferred.

Relative to the reviewer-requested inventory, five transformation classes were tested directly, three had only partial proxies, and three were untested. In particular, arbitrary prefix/suffix insertion, general case conversion, and general punctuation edits were not tested; line wrapping, email/markdown quoting, and sentence-boundary changes were only partially represented. First-divergence labels locate the first observed mismatch among tokenization, context, rank, or payload stages, but are descriptive and cannot by themselves establish why a transformation failed (Supplementary Fig. S1 and Supplementary Tables S7 and S8).

\begin{figure}[H]
\centering
\includegraphics[width=0.97\textwidth]{figures/robustness_recovery_compact.pdf}
\caption{\label{fig:robustness}\textbf{Exact-copy fragility.} Recovery proportions and Wilson 95\% intervals are shown for the principal replay and transmission conditions. Saved token IDs preserve the strongest replay contract. Visible-text detokenization and retokenization already reduce recovery. NFKC normalization, quote conversion, and trimming reduce it further, while markdown copy/paste, paraphrase, and cross-model decoding yield no successful payloads. Final 10\% tail truncation removes only declared non-payload tail tokens and must not be interpreted as arbitrary truncation tolerance. The full transformation matrix, denominators, limited canonicalization cells, and first-divergence diagnostics appear in Supplementary Figure S1 and Supplementary Note S4.}
\end{figure}

\subsection*{Forced span versus full message}\label{sec:results-quality}

Across 1,440 paired segmented trials, overall forced-span mean token log probability was -5.090904 and overall full-message mean token log probability was -1.424774, for an overall paired gain of +3.666130 (Fig.~\ref{fig:forced-full}). Each of the six model and topic schedule cells contained 240 paired trials, and all six 2,000-resample paired payload-bootstrap confidence intervals for the gain excluded zero.

Natural tails contain no payload ranks and add no payload capacity. The higher full-message average therefore reflects the inclusion of non-payload greedy continuations and must not be interpreted as improved quality on the payload-bearing forced span. Automated readability and surface diagnostics remain in Supplementary Fig. S2 and Note S5; they are not human evaluation.

\begin{figure}[H]
\centering
\includegraphics[width=0.97\textwidth]{figures/forced_span_full_message_quality.pdf}
\caption{\label{fig:forced-full}\textbf{Forced span versus full message.} Mean token log probability is paired within 1,440 segmented trials across three models and single-topic or multi-topic schedules, with 240 trials per model schedule cell. Error bars are 2,000-resample paired payload-bootstrap 95\% confidence intervals; all six intervals for the full-minus-forced gain exclude zero. Natural tails contain no payload ranks and add no payload capacity, so the improved full-message average is not an improvement to the payload-bearing forced span.}
\end{figure}

\subsection*{Ablations, dynamic stopping, and prompt-topic effects}\label{sec:results-ablations}

The locked compact ablation extraction compared 48 paired payload groups across three models and is exploratory/post-outcome (Table~\ref{tab:reviewer-analyses}). A 32-token lead-in, relative to no lead-in, reduced mean log probability by 1.652327 (95\% CI \([-1.749091,\allowbreak -1.553103]\)) and added 159.590278 full-message tokens (CI \([125.829861,\allowbreak 191.030035]\)). This adverse result is consistent with added context and boundary cost; the contrast does not prove a single failure mechanism. Increasing the segment size from 8 to 32 ranks increased effective rate by 1.231461 bits per full token (CI \([1.073116,\allowbreak 1.388188]\)) and reduced full length by 174.500000 tokens (CI \([-208.405035,\allowbreak -140.573264]\)), trading fewer tails for longer forced spans.
The lead-in and segment results locate different costs. A lead-in precedes the forced span, consumes visible tokens, and changes the distribution used for every encoded rank; it therefore cannot be treated as a harmless stylistic prefix. Larger segments keep the forced count fixed but require fewer prompt resets and tail opportunities. Their higher effective rate is consistent with reduced overhead, while the forced span itself lasts longer before an ordinary greedy continuation begins.

Removing the safe-text filter changed mean log probability by approximately -0.0018 (canonical CSV estimate -0.001797; CI \([-0.037738, 0.035394]\)) and effective rate by -0.062969 bits per full token (CI \([-0.134265, 0.007651]\)); both intervals included zero. The Mistral roundtrip-stable-filter cell was unavailable for 48 work units, although the no-filter contrast included all three models. The filter results therefore do not establish a quality benefit.
The near-zero selected filter estimates are important null evidence. The deterministic filter remains necessary to define an agreed candidate set for the canonical protocol, but this ablation did not show that it improved the selected likelihood or rate endpoints. Nor can the unavailable stable-filter Mistral condition be treated as evidence for or against that condition and it is an absence caused by the frozen isolated-vocabulary criterion.
Compared with dynamic stopping, no tail increased effective rate by 3.171551 bits per full token and shortened messages by 254.951389 tokens. This arithmetic consequence of removing non-payload text is not evidence of better cover quality. A fixed sentence-tail policy of \(20\) to \(60\) tokens shortened messages by 65.708333 tokens relative to dynamic stopping, whereas its selected effective-rate interval included zero. Dynamic stopping remains a declared heuristic with an emergency cap, not a semantic judgment by readers.

The tail ablation also prevents a misleading optimization conclusion. Maximizing bits per full token alone selects no tail because every added non-payload token enlarges the denominator. RankCloak introduced tails for a different, unverified purpose: allowing a forced fragment to continue to a rule-defined stopping point. Without human ratings, the experiment measures the length/rate price of that choice but cannot determine whether the price purchases perceived naturalness.
Across 15 locked topic contrasts, one Llama artifact-count comparison excluded the null; multi-topic versus single-topic expected count ratio 1.276361 (95\% CI \([1.114743, 1.461409]\); Holm-adjusted \(p=0.003294\)). The other 14 adjusted \(p\)-values were at least 0.396248. This compact extraction was performed post-outcome without refitting models and is exploratory as it does not support a general prompt-topic advantage.
The solitary non-null row concerns one model and one artifact-count outcome. It did not recur across the other model/outcome combinations after adjustment. Accordingly, prompt diversity is retained as design coverage rather than promoted as a performance-enhancing intervention. Reviewer-relevant ablations, exploratory tail-gain subgroups, and all topic contrasts appear in Supplementary Figs. S3 and S4 and Supplementary Tables S10 and S11; the complete 60-row contrast table remains machine-readable in the public repository.

\begin{table}[H]
\centering
\scriptsize
\setlength{\tabcolsep}{2.5pt}
\begin{tabularx}{\textwidth}{L{0.13\textwidth}L{0.13\textwidth}Y Y L{0.105\textwidth}Y}
\toprule
Analysis & Sample/group & Principal estimate & Uncertainty or adjusted test & Classification & Main limitation \\
\midrule
32-token lead-in vs none & 48 paired payloads; 3 models & Mean log probability -1.652327; full length +159.590278 & CIs \([-1.749091, -1.553103]\); \([125.829861, 191.030035]\) & Exploratory, post-outcome & Boundary/context interpretation is not causal \\
Segments of 32 versus 8 ranks & Same & Effective rate +1.231461 bits/full token; length -174.500000 & CIs \([1.073116, 1.388188]\); \([-208.405035, -140.573264]\) & Exploratory, post-outcome & Fewer tails accompany longer forced spans \\
No filter vs safe filter & Same & Log probability -0.001797; rate -0.062969 & CIs \([-0.037738, 0.035394]\); \([-0.134265, 0.007651]\) & Exploratory, post-outcome & Both selected intervals include zero; stable-filter Mistral cell unavailable \\
Tail policy & Same & No tail: rate +3.171551, length -254.951389; fixed tail: length -65.708333 & Selected 95\% CIs in Supplementary Note S6 & Exploratory, post-outcome & Rate gain from deleting non-payload text is not quality gain \\
Topic schedule & 15 locked contrasts & Llama artifact-count ratio 1.276361 & CI \([1.114743, 1.461409]\); Holm \(p=0.003294\); other \(p_{\rm Holm}\geq0.396248\) & Exploratory, post-outcome & One selected model/outcome contrast excludes null \\
Generation overhead & 18 model and protocol cells; 12,144 matched rows & Mean \([1.674368, 16.087628]\) s; \([13.631855, 100.851772]\) bit/s & Payload-bootstrap intervals in Supplementary Note S9 & Supporting measurement & Inclusive wall time; no CPU time or warm-up repeats \\
\bottomrule
\end{tabularx}
\caption{\label{tab:reviewer-analyses}\textbf{Compact reviewer-requested analysis and operational summary.} Differences are level minus canonical unless otherwise stated. Ablation and topic rows are a locked exploratory/post-outcome extraction from retained analyses; no new model was fit. Reviewer-relevant ablations, complete topic rows, availability, and timing cells are in the Supplementary Information; the full 60-row contrast table remains machine-readable in the repository.}
\end{table}

\subsection*{Neural steganalysis}\label{sec:results-detector}

\DIFdelbegin \DIFdel{Both classifiers strongly separated RankCloak covers from matched controls (Fig.~\ref{fig:detector}). TextCNN matched }\DIFdelend \DIFaddbegin \DIFadd{On the rebuilt strictly deduplicated matched test, TextCNN }\DIFaddend ROC-AUC was \DIFdelbegin \DIFdel{0.990190 (payload-group bootstrap }\DIFdelend \DIFaddbegin \DIFadd{0.9868 (}\DIFaddend 95\% CI \DIFdelbegin \DIFdel{\([0.987740,\allowbreak 0.992682]\)) with balanced accuracy 0.951010 (CI \([0.943456,\allowbreak 0.958466]\)). }\DIFdelend \DIFaddbegin \DIFadd{\([0.9819,0.9913]\)), }\DIFaddend DeBERTa-v3-base \DIFdelbegin \DIFdel{ROC-AUC was 0.999834 (CI \([0.999696,\allowbreak 0.999939]\))with balanced accuracy 0.991919 (CI \([0.987654,\allowbreak 0.995665]\)). High performance is adverse evidence for concealment, not a security feature of RankCloak. This test directly challenges the premise that removing visible artifact syntax is sufficient for concealment. The corpus}\DIFdelend was \DIFdelbegin \DIFdel{balanced between RankCloak and matched ordinary controls and split by payload group, so the matched estimates cannot be attributed to class imbalance or row-level payload leakage. Both a compact convolutional model and a pretrained transformer found strong signal.
The evidence is limited to this constructed corpus and training procedure, but within that scope it contradicts any claim of practical undetectability.
}\DIFdelend \DIFaddbegin \DIFadd{0.9994 (CI \([0.9984,1.0000]\)), and the exact-model-aware surprisal attacker was 0.999995 (CI \([0.999970,1.000000]\); Fig.~\ref{fig:detector}). Their normalized empirical 0--1\% partial AUCs were 0.6652, 0.9948, and 0.999460. These values show that the encoding is often detectable; classification of a cover does not reveal the concealed artifact or restore its recognizable surface syntax.
}\DIFaddend

\DIFdelbegin \DIFdel{Generalization depended on architecture and split. Across 18 held-out templates, TextCNN }\DIFdelend \DIFaddbegin \DIFadd{Leave-one-family-out }\DIFaddend ROC-AUC \DIFdelbegin \DIFdel{range was \([0.720421,\allowbreak 0.981705]\) and its balanced-accuracy range was \([0.565909,\allowbreak 0.919318]\); DeBERTa ranges were \([0.996291,\allowbreak 0.999979]\) and \([0.952273,\allowbreak 0.995455]\). Leave-one-model }\DIFdelend \DIFaddbegin \DIFadd{ranged from 0.9186 to 0.9784 for TextCNN, 0.9972 to 1.0000 for DeBERTa, and 0.999943 to 1.000000 for the model-aware detector. These fits exclude the target family from training and validation and test unseen-family transfer within the three-family corpus, not arbitrary future models. At validation-frozen nominal 1\% FPR thresholds, matched test TPR/observed FPR was 0.8317/0.0127 (19/1,491 false positives), 0.9913/0.0007 (1/1,491), and 0.9987/0.0020 (3/1,491), respectively. The matched validation and test sizes support empirical 0.1\% reporting; at that target the corresponding TPRs were 0.4829, 0.9886, and 0.9987, with 2, 1, and 3 observed false positives. Family holdouts had only 497 test negatives and did not support 0.1\% estimates; at frozen 1\% thresholds their observed false-positive counts reached 35, 28, and 3, showing calibration drift.
}

\begin{figure}[H]
\centering
\DIFaddFL{\includegraphics[width=0.97\textwidth]{figures/detector_model_family_generalization.pdf}
}\caption{\label{fig:detector}\textbf{\DIFaddFL{Strictly deduplicated detector evaluation.}} \DIFaddFL{Matched and leave-one-family-out ROC-AUC results are shown for two data-adaptive text-only detectors and the exact-model-aware surprisal attacker. Higher performance indicates that encoded covers can be distinguished from matched controls. Components connected by declared exact or near-duplicate relations do not cross partitions; the rule does not detect semantic paraphrases below its threshold. Supplementary Note S12 reports full split counts, validation-frozen low-FPR operating points, exact false-positive counts, and unavailable 0.1\% conditions.}}
\end{figure}

\subsection*{\DIFadd{Human controls, entropy gating, and quantization}}\label{sec:results-v3}

\DIFadd{The bounded human-authored secondary analysis selected 250 Dolly responses across 16 available templates, with no exact or declared near-duplicate links within that corpus or to the primary corpus. Against these negatives, }\DIFaddend ROC-AUC \DIFdelbegin \DIFdel{ranges were \([0.925623,\allowbreak 0.973650]\) }\DIFdelend \DIFaddbegin \DIFadd{was 0.7537 }\DIFaddend for TextCNN and \DIFdelbegin \DIFdel{\([0.996772,\allowbreak 0.999979]\) for DeBERTa; leave-one-codec ranges were \([0.966532,\allowbreak 0.998634]\) }\DIFdelend \DIFaddbegin \DIFadd{0.9846 for DeBERTa. At thresholds selected only on matched ordinary-model validation controls, 152/250 }\DIFaddend and \DIFdelbegin \DIFdel{\([0.997917,\allowbreak 1.000000]\). These are ranges across distinct held-out splits, not confidence intervals for one pooled estimate. The TextCNN template range shows meaningful distribution sensitivity ; a high matched score alone does not specify behavior for every new prompt. DeBERTa remained near perfect over the evaluated held-out regimes, which strengthens the adverse conclusion for these model, prompt, and codec families. It still does not establish performance on arbitrary human text, unseen generation systems, or future RankCloak modifications. The ranges deliberately preserve heterogeneity instead of averaging unlike test sets}\DIFdelend \DIFaddbegin \DIFadd{97/250 human texts were false positives. These high rates reveal detector-domain sensitivity and may reflect generic machine-versus-human differences; the analysis is not a participant study, human naturalness rating, or evidence that readers cannot notice RankCloak covers}\DIFaddend .

\DIFdelbegin \DIFdel{Exact text-identity }\DIFdelend \DIFaddbegin \DIFadd{Entropy gating produced a capacity--length--detectability tradeoff (Table~\ref{tab:v3-summary}). From ungated through median and strict gates, mean fixed-payload capacity conditional on completion declined from 7.4001 to 5.7920 and 4.0513 bits/token, while mean length ratio increased from 1.0000 to 1.3743 }\DIFaddend and \DIFdelbegin \DIFdel{token-sequence leakage checks passed for all 28 frozen splits. A stricter lexical diagnostic nevertheless found 53 normalized near-duplicate pairs at the declared threshold, including 44 cross-payload pairs, affecting 16 splits. Excluding affected test payload groups without retraining changed ROC-AUC by -0.001925 to +0.004568 for TextCNN and -0.000038 to +0.000166 for DeBERTa; balanced-accuracy changes were similarly small. This sensitivity does not remove the design limitation because it is diagnostic, threshold-dependent, and does not rebuild the training sets.
Fixed-seed repetition on the same CUDA device reproduced scores exactly for both architectures. It supplies same-device repeatability only; no CPU}\DIFdelend \DIFaddbegin \DIFadd{2.7615. The strict condition completed 114}\DIFaddend /\DIFdelbegin \DIFdel{GPU equivalence experiment was completed. Small exclusion deltas show that the observed metrics were not numerically overturned by this particular thresholded removal, but they cannot prove independence of lexical structure.
The sensitivity discarded affected test payload groups after training and did not reconstruct all split assignments. Supplementary precision,Brier score,}\DIFdelend \DIFaddbegin \DIFadd{120 payloads; its six maximum-budget failures are retained outcomes, and its capacity denominator is the 114 completed cases. Text-only ROC-AUC declined under the strict gate to 0.8850 }\DIFaddend and \DIFdelbegin \DIFdel{low-false-positive sensitivity are likewise labeled exploratory. Same-device exact repetition verifies the saved CUDA procedure, not portability across processor types (Supplementary Fig. S5 and Supplementary Tables S12 and S13).
}\DIFdelend \DIFaddbegin \DIFadd{0.8893, but exact-model-aware ROC-AUC remained 0.9904. Visible-text recovery was 62/120, 65/120, and 68/120 across the three conditions, so gating neither established undetectability nor reliable visible-text transport.
}\DIFaddend

\DIFdelbegin %DIFDELCMD < \begin{figure}[H]
%DIFDELCMD < %%%
\DIFdelendFL \DIFaddbeginFL \DIFaddFL{For the single pinned Qwen revision, saved-ID recovery was 960/960 for both Q4\_K\_M and Q8\_0, whereas visible-text recovery was 732/960 (76.25\%) and 758/960 (78.96\%). The paired Q8-minus-Q4 difference was 2.71 percentage points (grouped-bootstrap 95\% CI \([-0.73,6.35]\)); paired outcomes were 588 both, 144 Q4 only, 170 Q8 only, and 58 neither, with exact two-sided McNemar \(p=0.1582\). Thus there was no statistically supported Q8 recovery advantage. Independently generated paths diverged substantially (mean positionwise match 0.0155), while cross-quantization detection remained very strong in both directions.
}

\begin{table}[H]
\DIFaddendFL \centering
\DIFdelbeginFL \DIFdelFL{\includegraphics[width=0.97\textwidth]{figures/neural_detector_performance_compact.pdf}
}\DIFdelendFL \DIFaddbeginFL \footnotesize
\setlength{\tabcolsep}{3pt}
\begin{tabularx}{\textwidth}{L{0.17\textwidth}L{0.25\textwidth}L{0.24\textwidth}Y}
\toprule
\DIFaddFL{Analysis or condition }& \DIFaddFL{Recovery or completion }& \DIFaddFL{Capacity and length }& \DIFaddFL{Detection or paired inference }\\
\midrule
\DIFaddFL{Strict detector split }& \DIFaddFL{15,838 retained rows; 9,904/2,952/2,982 train/validation/test }& \DIFaddFL{Deduplication before features; zero declared cross-partition leakage }& \DIFaddFL{Matched AUC 0.9868/0.9994/0.999995 }\\
\DIFaddFL{Human-authored controls }& \DIFaddFL{250 frozen test negatives; no human labels used for fitting }& \DIFaddFL{Length-matched within 16 available templates }& \DIFaddFL{TextCNN/DeBERTa AUC 0.7537/0.9846; false positives 152/250 and 97/250 }\\
\DIFaddFL{Entropy ungated }& \DIFaddFL{Complete/saved ID 120/120; visible 62/120 }& \DIFaddFL{7.4001 bit/token, \(n=120\), conditional on completion; length ratio 1.0000 }& \DIFaddFL{AUC TextCNN/DeBERTa/model-aware 0.9953/0.9999/1.0000 }\\
\DIFaddFL{Entropy median gate }& \DIFaddFL{Complete/saved ID 120/120; visible 65/120 }& \DIFaddFL{5.7920 bit/token, \(n=120\), conditional on completion; length ratio 1.3743 }& \DIFaddFL{AUC 0.9825/0.9785/0.9999 }\\
\DIFaddFL{Entropy strict gate }& \DIFaddFL{Complete/saved ID 114/120; visible 68/120; six budget failures }& \DIFaddFL{4.0513 bit/token, \(n=114\), conditional on completion; length ratio 2.7615 }& \DIFaddFL{AUC 0.8850/0.8893/0.9904 }\\
\DIFaddFL{Q4\_K\_M versus Q8\_0 }& \DIFaddFL{Saved ID 960/960 each; visible 732/960 versus 758/960 }& \DIFaddFL{Independent-path position match 0.0155 }& \DIFaddFL{Difference 2.71 points, CI \([-0.73,6.35]\); McNemar \(p=0.1582\) }\\
\DIFaddFL{Cross-quantization detection }& \DIFaddFL{Q4-to-Q8 and Q8-to-Q4 held-payload tests }& \DIFaddFL{One Qwen revision; two quantizations }& \DIFaddFL{AUC TextCNN 0.9927/0.9944; DeBERTa and model-aware 1.0000/1.0000 }\\
\bottomrule
\end{tabularx}
\DIFaddendFL \caption{\DIFdelbeginFL %DIFDELCMD < \label{fig:detector}%%%
\textbf{\DIFdelFL{Neural detectability.}} %DIFAUXCMD
\DIFdelFL{Matched ROC-AUC }\DIFdelendFL \DIFaddbeginFL \label{tab:v3-summary}\textbf{\DIFaddFL{Compact summary of the focused V3 analyses.}} \DIFaddFL{Entropy capacity is fixed-payload bits per generated token }\DIFaddendFL and \DIFdelbeginFL \DIFdelFL{balanced accuracy are shown with payload-group bootstrap 95\% confidence intervals}\DIFdelendFL \DIFaddbeginFL \DIFaddFL{is conditional on completion}\DIFaddendFL ; \DIFdelbeginFL \DIFdelFL{held-out-template}\DIFdelendFL \DIFaddbeginFL \DIFaddFL{entropy recovery proportions use all 120 attempts. Detector triples are TextCNN}\DIFaddendFL , \DIFdelbeginFL \DIFdelFL{leave-one-model}\DIFdelendFL \DIFaddbeginFL \DIFaddFL{DeBERTa-v3-base}\DIFaddendFL , and \DIFdelbeginFL \DIFdelFL{leave-one-codec results are summarized as ranges across frozen splits}\DIFdelendFL \DIFaddbeginFL \DIFaddFL{exact-model-aware surprisal}\DIFaddendFL . \DIFdelbeginFL \DIFdelFL{Higher performance is adverse }\DIFdelendFL \DIFaddbeginFL \DIFaddFL{Empirical 0.1\% operating points are unavailable }\DIFaddendFL for \DIFdelbeginFL \DIFdelFL{RankCloak concealment }\DIFdelendFL \DIFaddbeginFL \DIFaddFL{the 250-negative human comparison and every entropy and quantization subgroup }\DIFaddendFL because \DIFdelbeginFL \DIFdelFL{it indicates that covers are distinguishable from matched controls. Both classifiers completed all 28 regimes. Exact identity leakage checks passed}\DIFdelendFL \DIFaddbeginFL \DIFaddFL{each has fewer than 1}\DIFaddendFL ,\DIFdelbeginFL \DIFdelFL{but a declared lexical near-duplicate diagnostic affected 16 splits; exclusion sensitivity produced small observed changes without eliminating that limitation}\DIFdelendFL \DIFaddbeginFL \DIFaddFL{000 negatives}\DIFaddendFL . \DIFdelbeginFL \DIFdelFL{The full six-panel view and exploratory precision, Brier}\DIFdelendFL \DIFaddbeginFL \DIFaddFL{Detailed denominators}\DIFaddendFL , \DIFdelbeginFL \DIFdelFL{low-false-positive}\DIFdelendFL \DIFaddbeginFL \DIFaddFL{intervals}\DIFaddendFL , \DIFdelbeginFL \DIFdelFL{leakage}\DIFdelendFL \DIFaddbeginFL \DIFaddFL{thresholds}\DIFaddendFL , and \DIFdelbeginFL \DIFdelFL{repeatability }\DIFdelendFL \DIFaddbeginFL \DIFaddFL{condition-level }\DIFaddendFL results appear in Supplementary \DIFdelbeginFL \DIFdelFL{Figure S5}\DIFdelendFL \DIFaddbeginFL \DIFaddFL{Notes S11--S15}\DIFaddendFL .}
\DIFdelbeginFL %DIFDELCMD < \end{figure}
%DIFDELCMD < %%%
\DIFdelend \DIFaddbegin \end{table}
\DIFaddend

\subsection*{Computational overhead and generalizability}\label{sec:results-overhead}

Across the 18 primary model and protocol cells, mean inclusive generation time ranged from 1.674368 to 16.087628 seconds. Mean payload throughput ranged from 13.631855 to 100.851772 bits/s, and effective rate from 0.710004 to 7.322003 bits per full token. These saved wall times include wrapper setup and decoding scope defined by the harness; CPU time and repeated warm-up measurements were not retained.

The approximately order-of-magnitude spread across cells shows that representation and model choice materially affect operational cost in the retained environment. Compact direct ranks can reduce forced work, whereas byte expansion and repeated segmented continuations increase it. Throughput and effective rate answer different questions: the first normalizes payload bits by wall time, while the second normalizes by visible full-message tokens. Neither is a deployment benchmark outside the recorded wrapper and device.
On the fixed CUDA benchmark, representative TextCNN wall/training times were 91.358/9.498 seconds with 306.25 MB peak allocated VRAM; DeBERTa-v3-base times were 1,536.661/1,425.063 seconds with 6.38 GB. Those device-specific measurements do not generalize across hardware. Evaluation across three model families and exact-copy Spanish and Simplified Chinese trials broadens the controlled design.% It does not validate very-large proprietary models, broad multilingual cover quality, operational channels, or real-world deployment.

The detector timings also illustrate the cost of the stronger adverse test: the pretrained architecture required substantially more saved time and memory than TextCNN in the single retained benchmark. CPU neural training was infeasible and was not used as a comparison point. These measurements support transparent study accounting, not a claim that the approach is efficient enough for an operational sender, receiver, or monitor (Supplementary Fig. S7 and Supplementary Tables S15 and S16).




\section*{Discussion}\label{sec:discussion}

RankCloak provides a reproducible method for converting serialized cryptographic artifacts into language-model next-token ranks and \DIFdelbegin \DIFdel{measuring the cover textthat realizes them}\DIFdelend \DIFaddbegin \DIFadd{realizing those ranks as cover text}\DIFaddend . Its representation choices expose a central \DIFdelbegin \DIFdel{trade-off}\DIFdelend \DIFaddbegin \DIFadd{tradeoff}\DIFaddend : direct subwords can be compact but leave \DIFdelbegin \DIFdel{the }\DIFdelend rank pressure uncontrolled, whereas ASCII-byte and hex-nibble codecs bound the permitted rank by spending more forced positions. Segmentation redistributes those positions across contexts, and natural tails add non-payload length. The formal and empirical capacity relationships make these costs explicit rather than treating longer cover text as additional payload capacity.
\DIFdelbegin \DIFdel{Complete recovery in the evaluated primary and secondary-language matrices establishes deterministic codec and replay fidelity under the saved-token exact-copy contract. It does not establish an ordinary text channel. Detokenized visible text could change tokenization without an apparent edit, and common normalization, formatting, rewriting, or model-mismatch conditions sharply reduced or eliminated recovery . The contrast between perfect }\DIFdelend \DIFaddbegin

\DIFadd{The principal demonstrated transformation is that the recognizable surface representation of a synthetic cryptographic artifact can be absent from the generated cover while an authorized receiver reconstructs it under the supported shared configuration. In the }\DIFaddend saved-token \DIFdelbegin \DIFdel{replay and fragile }\DIFdelend \DIFaddbegin \DIFadd{condition, retaining the token stream isolates tokenization fidelity, but recovery still requires contextual rank reconstruction and codec inversion; ordinary }\DIFaddend visible-text \DIFdelbegin \DIFdel{transmission is therefore not a secondary implementation detail; it defines the boundary of the recovery claim}\DIFdelend \DIFaddbegin \DIFadd{transfer is the stricter and materially less reliable condition. Steganalysis addresses another endpoint: a detector can often flag the presence of encoding, weakening covert-channel stealth, without thereby observing the original artifact or recovering its serialization}\DIFaddend .

Cover-quality evidence remained protocol-dependent. Full-message mean log probability was higher than forced-span mean across paired segmented trials, but the added natural tails carry no ranks and the gain does not improve the payload-bearing positions or add capacity. Automated readability, surface diagnostics, and model-based contrasts remain computational proxies rather than human judgments, and no participant evidence supports a claim of human-perceived naturalness. \DIFdelbegin \DIFdel{Supporting analyses reinforce this interpretation without identifying a single optimum. }\DIFdelend Larger segments improved efficiency by reducing repeated tail overhead, while lead-ins worsened likelihood and length in the selected contrasts; removing tails mechanically improved rate but did not test perceived quality.
\DIFdelbegin \DIFdel{These exploratory comparisons show that optimization depends on which cost is prioritized rather than establishing a universally preferable protocol. Evaluation across three model families, diverse prompts, and two secondary languages broadens the controlled setting but cannot be extrapolated to arbitrary models, languages, or deployed channels.
}\DIFdelend

\DIFdelbegin \DIFdel{Neural steganalysis supplies the clearest adverse finding. Both detector architectures achieved very high matchedperformance, and the pretrained detector remained highly effective across the evaluated held-out regimes. Although split sensitivity and lexical overlap limit extrapolation, the observed distinguishability directly argues against practical concealment in this corpus. Exact recovery cannot compensate for a cover distribution that trained classifiers can identify with near-perfect accuracy.
Despite these adverse results, the study contributes a controlled framework for separating representation, recovery assumptions, capacity}\DIFdelend \DIFaddbegin \DIFadd{The focused steganalysis revision strengthens the adverse evidence without changing that endpoint distinction. Strict pre-feature component deduplication found zero leakage under the declared exact and character-\(n\)-gram rules, yet matched}\DIFaddend , \DIFaddbegin \DIFadd{unseen-family, and exact-model-aware detectors remained strong. Validation-frozen low-FPR results show that high ranking performance does not ensure threshold calibration after family or quantization shift. The human-authored comparison is informative but secondary: its high false-positive rates expose control-domain sensitivity and may reflect general machine-versus-human cues.
}

\DIFadd{Entropy gating reduced text-only detectability in the tested matrix, but lower capacity, greater }\DIFaddend cover length, \DIFdelbegin \DIFdel{automated quality, transformation sensitivity, and detectability. The expanded payload, prompt, model, and secondary-language design shows how these quantities can be evaluated without equating model likelihood with reader judgment or deterministic replay with channel robustness.
RankCloak therefore establishes a measurement framework and a set of reproducible failure surfaces for rank-based hiding, but it does }\DIFdelend \DIFaddbegin \DIFadd{six strict-gate completion failures, and continued strong exact-model-aware detection prevent treating it as a uniformly superior protocol. Likewise, both tested quantizations preserved saved-ID inversion, but visible-text recovery remained weaker, independently generated token paths diverged, and the paired analysis did not support a Q8 recovery advantage. These studies identify measurable design tradeoffs rather than eliminating detectability or transport limitations.
}

\DIFadd{RankCloak therefore contributes an artifact-specific method, evaluation methodology, and empirical evidence about concealing and recovering synthetic cryptographic artifacts whose original syntax would be conspicuous. The experiments do }\DIFaddend not establish a \DIFdelbegin \DIFdel{secure, undetectable, robust, or deployment-ready covert communication system}\DIFdelend \DIFaddbegin \DIFadd{cryptographic secrecy proof, universal steganographic security, or robustness under every text transformation. They instead define when rank-based artifact recovery succeeds, when visible transport fails, and how current attackers detect the encoded covers}\DIFaddend .


\section*{Responsible use and limitations}\label{sec:limitations}

RankCloak is dual use and can encode sensitive-looking artifacts, but it does not \DIFdelbegin \DIFdel{itself }\DIFdelend provide encryption, authentication, confidentiality, or a \DIFdelbegin \DIFdel{cryptographic-security proof; any protected payload still requires }\DIFdelend \DIFaddbegin \DIFadd{cryptographic secrecy proof; protected payloads require }\DIFaddend an external cryptographic system and an explicit threat model. \DIFdelbegin \DIFdel{Exact recovery depends on saved-token replay, whereas }\DIFdelend \DIFaddbegin \DIFadd{The study does not claim universal undetectability. Saved-token inversion depends on the declared protocol configuration, while }\DIFaddend visible-text \DIFdelbegin \DIFdel{transmission and common transformations are fragile and the tested perturbations do not establish broad edit robustness . Human-perceived naturalness, broad multilingual behavior, real-world deployment, and performance across arbitrary future models or communication channels were not established; the automated }\DIFdelend \DIFaddbegin \DIFadd{recovery is weaker and markdown transfer, paraphrase, and cross-model decoding failed in the bounded robustness sample. The payloads are synthetic cryptographic artifacts, so the evidence does not establish hiding or recovery for arbitrary natural-language messages. Automated readability and likelihood }\DIFaddend measures are not participant ratings, \DIFdelbegin \DIFdel{and detector performance }\DIFdelend \DIFaddbegin \DIFadd{no participant study was performed, and the 250-response Dolly control is a bounded secondary corpus that may expose generic machine-versus-human differences.
}

\DIFadd{Strict deduplication }\DIFaddend is limited to the \DIFdelbegin \DIFdel{evaluated corpus. The evidence therefore supports RankCloak as a controlled research framework, not as a secure, undetectable, robust, or deployment-ready covert channel; detailed }\DIFdelend \DIFaddbegin \DIFadd{declared normalization and character-\(n\)-gram threshold and cannot rule out semantic paraphrases. The three-family holdouts, two learned text detectors, and exact-model-aware attacker are not exhaustive. Small family, human, entropy, and quantization test sets cannot resolve empirical 0.1\% operating points, and frozen thresholds can drift across family or quantization. The entropy experiment uses 120 attempts per gate and two templates; its strict gate left six payloads incomplete at the maximum budget, and its reported capacity is conditional on 114 completions. The quantization comparison covers one Qwen revision, one backend, and only Q4\_K\_M and Q8\_0. Broader models, tokenizers, quantizations, languages, transformations, communication systems, and adversaries remain untested; complete }\DIFaddend boundaries appear in Supplementary \DIFdelbegin \DIFdel{Table S17}\DIFdelend \DIFaddbegin \DIFadd{Notes S11--S15}\DIFaddend .


\section*{Data and Code availability}\label{sec:data-code}

The RankCloak source code and computational data \DIFdelbegin \DIFdel{supporting this study }\DIFdelend are available at the GitHub repository, \url{https://github.com/mantzaris/llm-rankcloak}\DIFdelbegin \DIFdel{, and are archived in Zenodo }\DIFdelend \DIFaddbegin \DIFadd{. The currently published Zenodo record is available }\DIFaddend at \url{https://doi.org/10.5281/zenodo.21987450}. \DIFdelbegin \DIFdel{The archive contains the public payload corpus, analysis inputs, detector predictions and metrics, source tables, figures, configurations, and provenance required to inspect the reported computational results}\DIFdelend \DIFaddbegin \DIFadd{Before submission, the final V3 revision will be synchronized to Zenodo, the archived version will be checked for the cited computational revision, and the DOI wording will be updated to match the deposited version; this statement does not claim that those final synchronization steps have already occurred}\DIFaddend .

\bibliography{references}

\section*{Acknowledgements}

NA

\section*{Author contributions}

A.V.M. alone worked on this paper and project.

\section*{Funding}

NA

\section*{Competing interests}

The author declares no competing interests.

\end{document}
